% Auto-generated by scripts/06_emit_structure.py
% Source: scans 074–107
% Section id: chapter/3
% Phase 3 deliverable. Footnotes (Phase 4), citations (Phase 5),
% and Wade-Giles → Pinyin (Phase 6) are not yet applied here.

\chapter{Deciphering the Inscriptions}
\label{ch:3}

% source: scan 074, printed 57
\section[Introduction]{\origsecnum{3.1}Introduction}
\label{sec:chapters-ch03:3.1}

Introduction\footnote[1]{For the famous story, first published by 董 in 1930, about Wang 乙-jung's sickness and 刘 E's discovery of inscriptions on the ``dragon bones'' from a Peking pharmacy, see Lefeuvre (1975), p. 17, n. 1; 李 Chi (1977), pp. 7-11. 刘 E himself made no mention of the incident in his introduction to 铁-阴ün.} to the Scholarship
Few events have excited the minds of modern Chinese scholars like the discovery
of the oracle bones. In the three-quarters of a century that has elapsed since 1899, when
Wang 乙-jung and 刘 E first appreciated the true nature of these ``dragon bone'' inscriptions, a large body of books and articles-at the last count, over 180 books and over 970
monographs has appeared “like bamboo shoots after spring rain.'' In this chapter, I
will first introduce the essential works and indicate how they may be used.\footnote[2]{岛 (1971), preface, p. 1; cf. 长 拼'üan (1967), p. 841, who refers to approximately 1,000 articles and over 180 books by about 300 scholars.}³\footnote[3]{The progress of oracle-bone scholarship (chiaku-hsüeh) has involved several interrelated stages: the identification of the true nature of the inscriptions; the identification of the site from which the inscribed bone fragments were being taken; the collection and publication of those fragments; the eventual conduct of scientific excavations in the 小'un region; and, above all, the gradual decipherment, still continuing, of the graphs. The history of this progress has been well summarized elsewhere. See, in particular, Kaizuka (1946), pp. 191-201; (1967), pp. 8-60; Jung Keng (1947); 胡侯-宣üan (1955), pp. 35-42; 陳 孟家 (1956), pp. 50-54; 董 (1965), pp. 46-118. For Western-language accounts, see Creel (1938), pp. 1-16; Lefeuvre (1975), which is fundamental; and 李 Chi (1977).}

\section[Published Sources]{\origsecnum{3.2}Published Sources}
\label{sec:chapters-ch03:3.2}

Approximately 107,000 oracle-bone pieces have now been assembled in over
eighty collections, generally in the form of rubbings.\footnote[4]{Matsumaru (1973), pp. 19-22, lists eightyeight collections (excluding those like 拼编, 朝村, and 通选, which duplicate material previously published) and indicates the style 57 ] U L U}4 A list of the major collections is
% source: scan 075, printed 58
given in Bibliography A, together with the abbreviations used to refer to them in this
book. Mechanically reproduced copies about whose accuracy there should generally be
no question (sec.\footnote[5]{Oracle bones are usually cited by a number which may be of two styles. The first and more traditional style (used by such collections as 践授, 徐-编, and 铁-阴ün) refers to the rubbing by citing the chüan (if any), the page number, and the number of the fragment, counting in Chinese style sequence (i.e., down the columns, right to left) and continuously from recto to verso; e.g., 钱-编 6.39.6 (S190.4) refers to the bottom rubbing in the right-hand column of p. 39b of ch. 6 of Yin-hsü shu-ch'i ch'ien-pien. In some cases the work is divided into parts 1 and 2 (shang and hsia F) (as in 侯编 and Hsüts'un) or into topical sections (as in 傅印), but in all these cases the last part of the citation is to the page number and to the number of the fragment on that page. The second style of citation refers to the oracle bone by the rubbing number printed by it; in collections using this method (e.g., 甲编, 易编, 人文, Menzies), the numbers run consecutively. It is always necessary to distinguish between the plate number and the rubbing number. In 拼编, the rubbing number is given in standard Chinese numerals at the lower left of the rubbing plate; it is the second number given at the start of the k'ao-shih. The plate number is given in documentary Chinese numerals at the top right corner of the rubbing plate; it is the first number given at the start of the k'ao-shih and is there preceded by the words t'u-pan. Thus, 拼编 548 (rubbing number), for example, is 拼编 plate 519. (The first edition of 岛 Sōrui [sec. 3.3.2] refers to the 拼编 plastrons by plate number, but the additional rejoinings recorded in the second edition are cited, inconsistently, by the 拼编 rubbing number.) Note that the 拼编 situation is complicated by the presence, under the plastron, of the rubbing numbers (in arabic numerals) and the registration numbers (in smaller arabic numerals) of the constituent 易编 fragments. (On registration numbers, see \ref{ch:1} (see ch. 1), n. 31.) A similar multiplication of numbers occurs in Chuiho. It is standard practice in all cases to cite the rubbing number (particularly since the plate number is ambiguous if more than one rubbing is reproduced on a plate). Other conventions may be noted. The notations cheng F. and fan refer to the front and back of a bone or shell (cf. \ref{ch:1} (see ch. 1), n. 46). In some collections (e.g., 人文, Menzies, 明候), the letters ``a,'' ``b,'' and ``c'' refer to inscriptions on the front, back, and (if a scapula) socket, respectively; ``S'' and ``B'' stand for shell and bone (cf. \ref{ch:1} (see ch. 1), n. 31); and an asterisk indicates if the graphs were filled with colored pigment (see \ref{ch:2} (see ch. 2), n. 145).} 5.5), these rubbings are the primary sources for most 商 historians.''
Many of these collections, especially those published more recently, have been printed
together with modern transcriptions. They frequently include commentaries about particular graphs and matters of interpretation. A scholar working with any collection of
rubbings should first determine whether or not a k'ao-shih or shakubun exists; these are
either published in the rear of the rubbings volume (as in 拼编), in a separate ts'e (as
in 乙-chu or 乙-村'un), or in a subsequent volume (as in 甲编 k'ao-shih or 人文 shakubun).
Table 9 lists the commentaries which exist for the major published collections.
of reproduction used in each. On the number of
fragments, see appendix 3, n. 2. On the provenance
and, when known, the location of the fragments
reproduced in these collections, see Jung Keng
(1947), pp. 15-30; 胡侯-宣üan (1955), pp. 13-
35; Lefeuvre \listitem{1975} and his work in progress,
tentatively titled ``Bibliographical Studies about
the Source Materials of Oracle-Bone Inscriptions,''
which also discusses the quality of various editions
of the same collection. Tamada (1972 et seq.)
provides punctuated and annotated versions of
the authors' prefaces to the various collections
together with Japanese translations.[\textasciicircum{}6 \textasciicircum{}6]: A few collections of rubbings are also provided with indexes, arranged either by oraclebone form (e.g., Pu-tz'u, pt. 3), 说文 form (e.g., 人文 sakuin), or modern-graph form (e.g., 崔-编, 通选). But these adjuncts have been rendered virtually superfluous by the appearance of CKWP (sec. 3.3.1) and 岛 Sōrui (sec. 3.3.2).

% source: scan 076, printed 59
\section[Reference Works]{\origsecnum{3.3}Reference Works}
\label{sec:chapters-ch03:3.3}
Over\footnote[7]{Kaizuka (1967, p. 204, estimates that we know about 3,000 graphs on the oracle bones. about 800 of which can be interpreted; 长 拼'üan (1967), p. 829, refers to 3,320 graphs, of which we can recognize about 1,000, with 500 to 700 more identifiable but in dispute. 胡 Houhsüan (1945), p. 5b, speculated, extravagantly in my view, that the full oracle-bone vocabulary, were we able to recover it, would have included 5,000 to 6,000 words, of which we would be able to recognize from 2,000 to over 3,000. No exactitude is possible in such counts (for other examples, see 长 春术 [1972], p. 283); nor is anything more than a general estimate necessary. See too n. 20.} three thousand 商 graphs have been identified on the oracle bones, and
about half of these may be translated or identified with varying degrees of certainty. The
dictionaries and concordances listed below are indispensable to the study of the inscriptions.

\subsection[Dictionaries]{\origsecnum{3.3.1}Dictionaries}
\label{sec:chapters-ch03:3.3.1}

Dictionaries\footnote[8]{中-kuo k'e-hsüeh-yüan k'ao-ku yen-chiu-so ‹† 國科學院考古研究所,甲-骨 wen-pien 甲骨文編, 考古-hsüeh chuan-k'an yi-chung ti-shih-ssu hao 考古學專刋乙種第十四號(Peking,1965). This work, prepared under the general editorship of Sun Hai-po is a revised edition of his 甲-骨 wen-pien (Peiping, 1934), which identified some 500 to 600 graphs, taken from eight collections (cf. Teng and Biggerstaff [1971], p. 46). Some of the 4,672 entries in CKWP probably refer to the same words and could be combined. The 1,723 entries in the main section of the dictionary consist of words which have been ``identified'' either in the sense that a 说文 graph can be assigned (in 941 cases) or that a modern graphic reconstruction and putative position within the 说文 system of classification can be provided. A ho-wen section contains 371 entries for graphs which frequently appear in conjunction, such as those used in ancestral titles. The appendix contains 2,949 entries of graphs which have not been deciphered or classified (CKWP, preface, p. 1). The forty collections consulted are listed with their abbreviations between the appendix and the modern-character index. An edition of this dictionary was reissued by 乙-wen in Taipei in 1975 under the title 蛟城 chia-ku wen-pien 校正甲骨文編,}
The standard dictionary is 甲-骨 wen-pien (abbreviated ChIP), which contains
a total of 4,672 entries taken from forty collections of published inscriptions. This needs
to be supplemented, on occasion, with 徐 chia-ku wen-pien.''
The entries in both dictionaries, like those in many works of oracle-bone scholarship,
are arranged according to the 540 radicals of 说文. Each entry consists of the 说文
graph (when known) as heading, followed by hand-copied reproductions of the oracle-bone characters it is thought to represent. Occasionally, there is brief discussion of the
graph's meaning, but generally the 说文 graph is the only “definition” provided. The

[\textasciicircum{}9]: 金 祥-heng, 徐 chia-ku wenpien (Taipei, 1959). This is a supplement to Sun Hai-po's 1934 dictionary (referred to in n. 8). Based on thirty-eight published collections, it contains 1,048 oracle-bone graphs identified with 说文 equivalents, 1,585 graphs not found in 说文, an appendix of ho-wen, and an appendix of some 300 unidentified graphs. There are two modern-character indexes for the graphs with and without 说文 equivalents. Cf. Noel Barnard in RBS 5 (1959), p. 463. All three dictionaries may need to be consulted on occasion. For example, for the graph yüanë, CKWP \listitem{1965} has more citations (thirty-two) than the 1934 edition (nine citations) and the 徐 chia-ku wen-pien (twelve citations) combined. But due to the decision of the 1965 editors (preface, p. 3) not to reproduce identical forms of the same character, the 1965 edition has many less entries for the graph chiu than the 1934 edition. 59 11 \} U 11 0
% source: scan 077, printed 60
modern-character indexes in the rear permit the reader to find a graph, provided he
already knows its modern form.
More useful as a dictionary is 甲-骨 wen-tzu chi-shih (abbreviated CKWT).\footnote[10]{李 小庭 *, 甲-骨 wen-tzu chishih X, 中-yang yen-chiu-yüan lishih yü-yen yen-chiu-so chuan-k'an chih wu-shih 中央研究院歷史語言研究所專利之五十,8 vols. (Nankang, 1965).} 10 Based
upon the graphs listed in 甲-骨 wen-pien \listitem{1934} and 徐 chia-ku wen-pien (1959), and also
organized according to the 说文 system, this handwritten work assembles in one place
the commentaries, sometimes abbreviated, of scholars who have discussed the etymologies
of individual oracle-bone graphs. At the end of each entry, 李 小庭 reviews the
evidence and gives his own conclusions. A table of contents and various indexes to the
main text, a supplement, and first appendix (for graphs not yet identified with certainty)
make the work convenient to consult. The opinions of some scholars have at times been
overlooked, but the work is an essential reference to early Chinese etymology, focusing
on earliest word usage as found in the 商 inscriptions.\footnote[11]{For more information about CKWT, see my account (forthcoming) in RBS 11. Recent studies of particular graphs may be found in 中-kuo wen-tzu X‡ (Taipei); there is an index in vol. 41 (1971). This periodical ceased publication, temporarily at least, as of vol. 52 (June 1974).} 11 It is not cited specifically on
many occasions in this book because I assume it to be a first source to which any scholar
will turn for information about particular graphs.\footnote[12]{The second edition of 岛 Sōrui (sec. 3.3.2) keys each entry to the proper page in CKWT and conveniently assembles these cross references at S602-607.} 12
甲-骨-hsüeh cannot be divorced from chin-shih-hsüeh.

商 bronze inscriptions which contain more than a minimal number of graphs are relatively few,
but since the language of the bronzes may throw light on the language of the bones, and
since our understanding of 商 graphs frequently depends on our understanding of
the graphs found on 周 bronzes, the scholar of 商 history should be prepared to
consult the two major bronze inscription dictionaries: 金-文-pien and 金-文 ku-lin.
金-文-pien (abbreviated CWP) contains in its main text entries for 1,894 graphs found in
商 and 周 bronze inscriptions, arranged according to the 说文 system.\footnote[13]{Jung Keng, 金-文-pien (Peking, 1959). This is a revised edition of Jung Keng, 金-文-pien (Tientsin [?], 1925). The 1959 edition has been reissued by Lien-kuan ch'u-pan she as the first part of 金-文 chenghsü-pien ho-ting-pen E✰★ (Taipei, 1971). For more information about the work, see Teng and Biggerstaff (1971), p. 47.}¹³ Each
entry is headed by the 说文 form or a modern graph reconstruction, followed by the
bronze forms listed in columns. A modern graph equivalent is listed under the first bronze
form with occasional brief discussion. The name of the vessel from which it comes is given
under every graph. An appendix with 1,199 entries gives graphs and insignia whose
modern forms have not been identified. This is followed by an index to the 3,165 vessels
consulted and an index of modern graphs. Like 甲-骨 wen-pien, 金-文-pien is less a
dictionary than a list of occurrences. Definitions are minimal.
金-文 ku-lin, which supersedes 金-文-pien, is a combined concordance and
% source: scan 078, printed 61
dictionary. 14 Following the order of the graphs in 金-文-pien (1959), it not only attempts
to reproduce every occurrence of the graph in question that may be found in the published
corpus of 商 and 周 vessels,\footnote[14]{周 Fa-kao with 长 Jih-sheng [Cheung Yat-shing], 徐 知-yi [Tsui Chee-yee] and Lin Chich-ming [Lam Kit-ming]林潔明,eds., 金-文 ku-lin 金文詰林, 16 vols. (Hong Kong, 1974-1975). For a review, see Mattos (1976).}\footnote[15]{Jung Keng's original corpus in 金-文pien \listitem{1959} has been expanded to include fiftythree more inscriptions from 郭 沫若 (1957).}\footnote[16]{Bernhard Karlgren, Grammata Serica Recensa (Stockholm, 1957).}\footnote[17]{Tōdō Akiyasu, Kanji gogen jiten (Tokyo, 1969). For a review of an earlier version of this work, see Pulleyblank (1968).}\footnote[18]{Katō Joken, Kanji no kigen 起原(Tokyo, 1970).}\footnote[19]{岛 Kunio, Inkyo bokuji sōrui , 2d. rev. ed. (Tokyo, 1971). The first edition (Tokyo, 1967), which is now less useful because it lacks the modern-character index, modern-character identifications at the entry heads, in-text cross references, cross references to CKIT, and 易编/拼编 identifications of the second edition, was reprinted in Taipei by T'aishun shu-chuin 1970. Despite the reduced page size of the second edition (1971), it is to be preferred in every respect to the first and is the edition usually cited in this book.}\footnote[20]{From the table of contents (S9-12) I estimate that the work contains about 3,095 entries. 岛, however, indicates that the 840 entries for which cross references are given to CKWT constitute no more than one-fifth of the entries (front matter, p. 3), so he would seem to place the number of entries at over 4,000. 61 ㄈㄈ U \} 0 } 15 it also quotes the sentence in which the graph appears.
Some twenty thousand of these inscription sentences are reproduced in the 1,894 entries.
In addition, commentaries by other scholars dealing with the graph or the sentence are
quoted in full, together with copious notations by 周 Fa-kao and his fellow editors.
Reproducing the inscriptional context and assembling the scholarship in this way, 金-文 ku-lin is a major reference work for etymological study.
The fact that these five dictionaries are all arranged by the 说文 system has
advantages. In the first place, the 540 radicals provide a convenient series of headings
under which to place graphs for which no modern form exists. Further, since the dictionaries are organized on the same principle, a scholar may move rapidly from one to
another; a graph at the start of chian 5 in one will be at the start of chüan 5 in all (it is for
this reason that old style citations are preferable to new [for example, CKWP, ch. 12,
p. 18b, rather than CKWP, p. 496]). The modern-character indexes at the rear permit
ready entry into the system for those not familiar with it.
Since the study of the inscriptions is not just the study of graphic forms but is also
the study of language (cf. sec. 3.5.2 on phonology), it goes without saying that a student
of the 商 inscriptions will consult other etymological dictionaries as the need arises.
Particularly useful are Grammata Serica Recensa (abbreviated GSR), 16 Kanji gogen jiten, 17
and Kanji no kigen. 18 These dictionaries attempt to establish phonetic and semantic word
families in early Chinese, and sometimes refer to the oracle inscriptions.

\subsection[Concordances and Compendiums]{\origsecnum{3.3.2}Concordances and Compendiums}
\label{sec:chapters-ch03:3.3.2}

Inkyo bokuji sōrui (abbreviated S) is a concordance of the oracle inscriptions in
sixty-three collections. 19 Over three thousand words have been classified. 20 Under each
% source: scan 079, printed 62
entry, 岛 Kunio has hand copied, in a standardized oracle-bone script (examples of
which grace the margins of this book) all the inscriptions in which the graph in question
appears, giving the citation for each. There are inevitable errors, inconsistent transcriptions, and omissions; some entries for the more common graphs have been abbreviated
(this is noted); and inscriptions published since 1959, together with some minor collections
published only in journals, are not included.\footnote[21]{It is true that the second edition identifies rejoined 易编 fragments that appear in the 拼编 volumes published down to 1967, i.e., down to 拼编 512 (for indexes to these rejoinings, see n. 72); but in terms of new fragments (as opposed to new pieces; cf. appendix 3, n. 1) no collection published since 1959 is included.} 21 But the work's fundamental contribution
is threefold. First, it attempts to provide a comprehensive listing of word usage in context.
Second, it provides control over word frequency (except, as already noted, in the case of
common graphs). Third. since 岛 copies the inscriptions under each entry in chronological order (cf.\ref{sec:chapters-ch04:4.1.1} (see  sec. 4.1.1)), it provides (within the limits imposed by the accuracy of his
transcription) a rapid introduction to the epigraphic evolution, if any, of individual
graphs (cf. sec. 4.3.1.4 and to their periodicity (appendix 5).
The work is also of major importance because, unlike the oracle-bone dictionaries
described above (sec. 3.3.1), it allows a graph entry to be found, even when the seal or modern
form of the graph is not known. This may be done by using the rear index (pp. 590—601),
arranged under the 164 oracle-bone ``radicals'' devised by 岛. Once the entry has
been found, researchers may study all occurrences of the graph in context; they may move
on to the commentaries which have been published for certain of the inscriptions cited;
or they may, in the 840 cases for which the cross reference is provided, move to the relevant
page in 李 小庭's 甲-骨 wen-tzu chi-shih.\footnote[22]{Examples of how this works are given in secs. 3.6.1 to 3.6.3; see too n. 12. For my fuller description and evaluation of Inkyo bokuji sōrui, see Keightley (1969a) and RBS 12 (forthcoming).}22 Inkyo bokuji sōrui and 甲-骨 wen-tzu
chi-shih are complementary reference works. Neither should be relied on as a primary
source for the inscriptions or commentaries. But they are invaluable guides to the sources
and scholarship, the two fundamental works in the library of any oracle-bone scholar.
Based on a corpus of fifty-eight published and seven unpublished collections, Yin-tai
chen-pu jen-wu t'ung-k'ao\footnote[23]{焦 Tsung-yi Yin-tai chen-pu jen-wu t'ung-k'ao, 2 vols. (Hongkong, 1959).} 23 prints in modern transcription all the divination charges associated
with the names of 137 diviners.\footnote[24]{Cf. \ref{ch:2} (see ch. 2), n. 13.} 24 Within each diviner's section, the inscriptions are subdivided by topic: rain, the ten-day week, eclipses, hunting, and so on. A compendium
rather than a concordance, 焦 Tsung-yi's work is less useful than 岛; first, because
it is organized by diviner and topic rather than by words; second, because divinations
lacking a diviner's name are not included; third, because his modern-character transcriptions, useful though they are to the beginner, introduce a considerable element of interpretation. But it is a valuable supplement, especially for studying the role of particular
diviners. Once the diviner's name and the topic of a particular charge are known, a modern
transcription may generally be found without difficulty.\footnote[25]{For a more detailed assessment of Yin-tai chen-pu jen-wu t'ung-k'ao, see Noel Barnard in Monumenta Serica 19 (1960), pp. 480-486. \_} 25
% source: scan 080, printed 62
Bibliographies

\subsection*{General Studies}

Four major bibliographies of oracle-bone studies have been published. The first
is 胡侯-宣üan's Wu-shih-nien chia-ku-hsieh lun-chu-mu, arranged topically and with
author and title indexes.\footnote[26]{胡侯-宣üan, Wu-shih-nien chiaku-hsüeh lun-chu-mu + (Shanghai, 1952; reissued, Hongkong, 1966). 胡 (1955), pp. 38-41, lists by category the principal books likely to be useful for beginners.} 26 Although useful, it has been superseded to a large extent by
Peng's ``An Annotated Bibliography of the Publications on the Inscribed Oracle Bones
Excavated from the Yin Ruins,\footnote[27]{Madge Shu-chi Peng, ``An Annotated Bibliography of the Publications on the Inscribed Oracle Bones Excavated from the Yin Ruins,'' Chinese Culture 6.3 (1965), pp. 97-149. The annotations are entirely bibliographical in nature.}'' 27 arranged by author, and by 董 and Huang's 徐
chia-ku nien-piao.\footnote[28]{董 作宾 and Huang Jan-wei , Hsu chia-ku nien-piao (Taipei, 1967). For further discussion of oracle-bone bibliographies, see my account (forthcoming) in RBS 12.} 28 This last is arranged chronologically, but the author index in the rear
permits it to be used as a bibliography. The most recent bibliography is Christian Deydier's
Les Jiaguwen: essai bibliographique et synthèse des études; this covers the period down to 1974
and provides romanizations for the works cited.\footnote[29]{Christian Deydier, Les Jiaguwen: essai bibliographique et synthèse des études, Publications de l'école française d'extrême-orient, vol. 106 (Paris, 1976).} 29 For current scholarship, section XIII.2
of the annual bibliography Tōyōgaku bunken ruimoku may be consulted.\footnote[30]{Kyōto daigaku jimbun kagaku kenyūjo Tōyōgaku bunken ruimoku * (Kyoto, yearly); title slightly different before 1963. Hayashi Minao attempted to provide annual bibliographies of oracle-bone and bronze inscription studies (Kōkotsugaku 6 [March 1958] to 9 [1961]); it is to be hoped that if Kōkotsugaku resumes publication, this bibliography will be continued. and \#8}30 Bibliographies
dealing with studies of bronze inscriptions should not be overlooked.\footnote[31]{Consult in particular Jung Yuan Jung Keng, 金-shih shu-lu mu (Shanghai, 1936), and its supplement (KK [1955.3], pp. 70-80) (reissued together as Chinshih shu-lu mu chi pu-pien [Taipei, 1971]); Hiroshima daigaku bungakubu Chugoku tetsugaku kenkyūshitsu ★\# 中國哲學研究室, Kimbunkankei bunken mokuroku 金文關係文獻目錄,mimeographed(Hiroshima, 1956). For more recent studies, the Toyogaku bunken ruimoku (n. 30) should be consulted.}31
General Studies
Several comprehensive studies of 商 culture, based in varying degree on the
information found on the oracle bones, may help the beginning student place the inscriptions in historical context. The best semipopular accounts are those by Itō Michiharu
(1967), Kaizuka Shigeki (1967), and Shirakawa Shizuka (1972); 李 Ya-nung \listitem{1955} is
also of interest. More technical studies, useful because they cite inscriptions copiously, at
times assuming the appearance of a catalogue raisonné, are those by Wu Tse (1953), 岛
Kunio (1958), Itō (1975, pt. 1), Akatsuka (1977), and, still the best comprehensive introduction, 陳 孟家 (1956). The study of 商 religion by 长 宗 \listitem{1970}
is also recommended since it translates nearly five hundred inscriptions.\footnote[32]{Consult too, the extensive review by Serruys (1974), which retranslates many of the same inscriptions. For some recent monographs on various aspects of 商 culture, see \ref{ch:5} (see ch. 5), n. 87. 63 U U CEE U 0} 32
% source: scan 081, printed 63

\section[Translation: General Considerations]{\origsecnum{3.5}Translation: General Considerations}
\label{sec:chapters-ch03:3.5}

Historians, grammarians, epigraphers, and linguists all have a contribution to
make to the translation of 商 inscriptions. Studies which place the inscriptions in a
cultural context may aid philologists and grammarians by establishing the frame of reference
in which the inscriptions were made. Philology, in Peter A. Boodberg's striking phrase, is
``an unending conversation with the dead.'' We cannot have that conversation without a
knowledge of their history, just as we cannot know their history without that conversation.
The whole issue of what is meant by the ``translation'' of a 商 inscription is a
thorny one.\footnote[33]{For a useful introduction to the problem, Smith (1972). see Serruys (1974), esp. pp. 19-21.} 33 Leaving aside the anthropological difficulties involved in trying to understand
one culture's conceptions from the viewpoint of another,\footnote[34]{On this issue, see in particular Winch \listitem{1964} and Lukes (1967), together with the other articles contained in Wilson (1970); Rorty (1972); J. Z.} 34 the technical problems alone
are intimidating. Phonetic considerations are of some assistance (sec. 3.5.2), but, as yet,
scholars must deal primarily with the written graphs, representations of a spoken language
we can no longer hear. To transcribe those 商 graphs as modern Chinese graphs (sec.
5.6), to say nothing of rendering those graphs into English, inevitably betrays the original
商 conceptions.
Given these limitations, the scholar should not merely leap, often in an act of semantic
faith, from the shape of the 商 graph to an apparent small-seal descendant and then
attempt to force the 说文 definition upon the 商. It is true that our ability to read
the 商 documents rests firmly on 周 and 汉 texts; if we did not have Western
周 bronze inscriptions, 周 texts, small-seal forms, and 说文 to assist us, it would
be virtually impossible to read the oracle-bone inscriptions. But the danger of contamination by later conceptions must always be evaluated before a translation of the 商
materials is attempted. As Britton noted long ago, ``tracing the stylization of a written
symbol through ten or twelve centuries is one thing, and determining its changing sense
from epoch to epoch is another.”\footnote[35]{Britton (1936), p. 207.} 35
This chronological and cultural distance between the 商 inscriptions and the first
Chinese texts that give us the subtle flavor and resonance of meaning on which true translation depends should make us cautious about coining 商 neologisms if the neologisms
depend upon classical overtones rather than 商 usage.\footnote[36]{On neologisms in Chinese translation, cf. Schafer (1972), pp. 150-151.} 36 We have no way other than
context of determining the contemporary flavor and emotional load of most of the 商
words with which we deal, a deficiency that makes our view of 商 history flat and
imprecise (cf.\ref{sec:chapters-ch05:5.2} (see  sec. 5.2)).
Any translation offered, therefore, should first take account of the various contextual
requirements which the word must satisfy in the 商 inscriptions. It is for this reason
that 岛 Inkyo bokuji sōrui (sec. 3.3.2) is valuable, since it provides us with copies of
% source: scan 082, printed 64
most of the published inscriptions which contain a particular graph or phrase. Ideally,
one must attempt to establish graphic, semantic, and, where possible, phonetic ties between
商 and later words, ties which show that no philological violence is being done by
taking the later word as the descendant of the 商 word. But it is not always possible
to make such filiations, and hence it is not always possible to make a precise translation
or to identify an equivalent modern Chinese word with confidence. In such cases, we
should not strive for an accuracy that may be misleading, but should rely, provisionally
at least, on a functional translation or kenning which on the basis of 商 usage appears
to carry the appropriate general meaning.\footnote[37]{The various ``disaster” words (e.g., [S209.2]; [S184.3]; [S303.1]) studied by Mickel \listitem{1976} typify the problem; it is difficult for us to distinguish them with appropriate translations in the way the 商 presumably did (cf. \ref{ch:2} (see ch. 2), n. 80). See too \ref{ch:2} (see ch. 2), n. 7, on the meaning of *tieng.} 37
The advance of oracle-bone scholarship depends to a large extent upon the improvement of previous translations.\footnote[38]{Strikingly divergent translations have been offered for the same inscription. In some cases the context is insufficient for any firm determination; for example, do the two barely discernible graphs shih hsiang (拼编 170.1) mean ``Promise ivory'' or ``Shoot arrows at elephants'' (Serruys [1974], p. 34)? Does 何編 153 (D) (S31.3) refer to the ``summer solstice'' or to the ``arrival'' of a person (董 [1954a], p. 127)? In other cases, the meaning of particular graphs is in dispute. To choose some examples virtually at random: (I) Does 侯编 2.40.15 (S412.3) mean ``Lady 京's millet will not be harvested'' (Serruys [1974], p. 35) or ``At Fu 京's millet planting, we will not perhaps offer libation'' (cf. sec. 3.6.3)? \listitem{2} Does 库-方 1649 (D) (S1.2) mean ``Do not make ascend and present as sacrifical victims 3,000 men'' (Serruys [1974], p. 70 and n. 59) or ``We should not raise (i.e., mobilize) 3,000 men'' (Keightley [1969], pp. 66-69)? \listitem{3} Does 徐-pu 1560 mean ``If we amputate (a criminal's) foot, will there be death'' (胡侯-宣üan [1973], p. 113) or (by analogy with the inscriptions at S12.4-13.3) ``Ling will not die”? \listitem{4} Does Houpien 2.9.1 (S362.3, but wrongly cited as 2.7.1) refer to the appearance of a great new star or to the linking of a great new earthwork to the mountains (see the various interpretations quoted by Ikeda [1964], 2.9.1)?} 38 Disagreements about translation usually involve the actual
wording (that is, what is to be seen on the bone), the meaning of the graphs, the grammar
of the inscription, or its punctuation. Ideally, all scholars who work with 商 inscriptions which involve major problems of interpretation should justify their translations with
notes or glosses.\footnote[39]{Good examples of glosses may be found in the commentaries to 拼编, 侯编 (Ikeda [1964]), and 人文. Glosses for particular words are conveniently assembled in CKWT (sec. 3.3.1).}39 Idiosyncratic and unsupported translations have no place in 商
studies.

\subsection[Grammar]{\origsecnum{3.5.1}Grammar}
\label{sec:chapters-ch03:3.5.1}

The rigorous study of oracle-bone grammar, depending upon a comprehensive
analysis of all the attested divination sentences in their full context, has hardly begun.
Scholars have assumed that the grammar of chia-ku-wen was substantially that of 周
dynasty wen-yen texts.\footnote[40]{E.g., 陳 孟家 (1956), pp. 13213465 D ]}40 This assumption is generally correct-verbs follow their subjects
and precede their objects, adjectives or adjectival clauses precede the nouns they modify,
% source: scan 083, printed 65
and so on-but we should be wary of automatically imposing wen-yen grammar on every
chia-ku text. First, because the oracle inscriptions were recorded several centuries prior to
the earliest 周 texts we possess and may involve more archaic grammatical forms;
second, because, being divination texts, they may manifest formulaic idiosyncrasies.\footnote[41]{The preverbal particle H, for example, is antecedent to the modal ch'i of later texts, but it has a characteristic function in the divination charges: it is associated with what are assumed to have been undesirable charges or those about whose desirability the diviner felt uncertain (cf. \ref{ch:2} (see ch. 2), n. 124), and it thus appears to have weakened the force of a charge; it may have served to indicate conditional subordination of the following clause. For the details of what is a complex argument, see Takashima (1973), pp. 267-305; Serruys (1974), pp. 24, 46-58. Takashima (1976), p. 24, however, argues that ``ch'i does not introduce subordinate clauses, and that the question of subordination must and can be deduced from context and sentence parallelism.'' I believe ch'i is best regarded as a particle of hypothesis and/or futurity (i.e., uncertain expectancy), and I translate it in the charges as ``perhaps'' (as in 拼编 8.2 [fig. 8]; 235.1 [fig. 11]; 247.2 [fig. 12]) and in prognostications either by ``if'' (拼编 247.1 [fig. 12]; 248.7 [fig. 13]) or by future uncertainty (拼编 19.9 [sec. 3.7.1.2]; 207.3 [sec. 3.7.4.2]; 京华 2 [sec. 2.8]).} 41
The standard introduction to oracle-bone grammar, despite its unsystematic and
dated nature, is still the chrestomathy of 陳 孟家, which leads the student through
word order, particles, time words, measures, pronouns, verbs, modifiers, numbers, demonstratives, connectives, prepositions, auxiliary verbs, negatives, omissions and abbreviations,
and sentence types.\footnote[42]{Ch’en 孟家 (1956), pp. 85–134.} 42 For more sophisticated studies indicating the linguistic methodology
to be employed and the kind of results to be gained, Takashima \listitem{1973} and Serruys \listitem{1974}
should be consulted.\footnote[43]{See too, Kryukov \listitem{1968} and (1973). 周 洪祥 \listitem{1969} provides much data, especially on the question of omissions. On measure words, see Huang 蔡-chün (1964). The studies of Huang 京-hsin \listitem{1958} and Kuan 谢楚 \listitem{1953} have been superseded. ``} 43 They deal with such matters as passives and causatives, conditional
constructions, word order and stress position, and the significance of the various negatives.\footnote[44]{Particular attention has been paid to wu `` (cf. n. 63), which seems to have been the strongest of the 商 negatives, with a sense of ``must not,' ``do not.'' It is likely, as Takashima (1973), p. 190, suggests, that in some cases wu had an obligatory sense; the frequent use of wu in the charge before verbs of sacrifice, for example, can be better translated as ``we need not (or should not) offer sacrifice'' rather than ``must not.'' It has also been proposed (Serruys [1974], pp. 44, 67; Takashima [1973], pp. 294-296, 399, n. 2) that wu be taken as a mark of subordination when it appears in the first of two clauses. Thus, 拼编 147.2, `` on the left side of the plastron, would be translated ``If we do not build a settlement, Ti will approve.'' And the presence of the wu would indicate that the matching, positive charge, 拼编 147.1\#4 on the right side of the plastron, should also be translated with a conditional first clause, ``If we perhaps build a settlement, Ti will not assist and approve.'' Similarly, 拼编 442.7, on the left side of the plastron, would be translated as ``On the next ting-mao, if we do not (perform a rain-seeking dance), it will not perhaps rain.'' This charge is matched by 并-编442.6,羽丁卯求舞ㄓ雨,on the right side of the plastron, ``On the next tingmao, if we perform a rain-seeking dance, there will be rain.'' In some cases, however, the ``if'' construction is less appropriate. To translate 拼编 12.2 (\ref{sec:chapters-ch03:3.7.1.2} (see see sec. 3.7.1.2)) as ``This season, if the king does not follow Wang 成 to attack the 夏 韦, (we) will not perhaps receive assistance in this case'' is unsatisfactory, for without the attack there would be no need for the assistance. In all cases where wu appears in the first of two clauses, therefore, I would suggest-and here I have benefited greatly from the advice of David Nivison (letter of 19 May 1976)-that the initial translation take wu in its obligatory sense, and that the positive, matching charge also be trans-} 44
It is to be hoped that work of this order will eventually lead to the writing of a modern
chia-ku-wen grammar.
% source: scan 084, printed 66

\subsection*{Phonology}
The inscriptions tell us so little about their sound that ``the problem of the pronunciation of 商 graphs'' has been declared ``nigh insurmountable.”\footnote[45]{Serruys (1974), p. 92, n. 6.}45 Since there is
no evidence of rhyme in the oracle inscriptions, the traditional and even modern approach
to their study has been resolutely graphemic. The systematic study of 商 phonology
has barely begun, and the inscriptions serve primarily as a testing ground for phonetic
reconstructions derived from Sino-Tibetan or Archaic Chinese. The extent to which such
reconstructions-on which no full agreement exists-may eventually assist our reading of
the inscriptions is not yet clear.\footnote[46]{Takashima (1973), appendix II, ``The Problem of Phonology,'' discusses the systems of phonetic reconstruction, the reasons for preferring that of 李 Fang-kuei \listitem{1971} to Karlgren's, the problem of defining archaic Chinese (ca. 700 B.C.) and whether or not it was phonetically uniform, and the need to discover ``Western 周” rhymes. Unless otherwise specified, I use, as a matter of convenience, the archaic reconstructions of Karlgren (all archaic reconstructions being preceded by an asterisk); GSR is readily available and conveniently arranged by word families. Tōdō \listitem{1969} may be consulted for suggested revisions of Karlgren's phonology and for its attempts to define the core meaning of the word families. 周 Fa-kao \listitem{1973} provides a useful tabulation of the archaic reconstructions proposed by 董 同何, Karlgren, and 周 Fa-kao. Pulleyblank will in the future provide a systematic account of his own reconstructions.}46
It must be remarked that the nature of the 商 script does not always facilitate our
studies. Certain graphs, at least, appear to have had no inherent phonetic, or even semantic,
value, but were used merely as shorthand notations, with the result that identical graphs
could have widely different meanings and sounds in different contexts. Documented cases
of such semantic and phonetic laxity are rare, but the possibility that one graph had two
or more radically different meanings and sounds must always give us pause.\footnote[47]{For example, the graph I might mean either *niam (modern jen E) or, after period II, *kung (modern kung I) (CKWP, \ref{ch:5} (see ch. 5), p. 2a; ch. 14, p. 14b); the graph 1 might mean either *nien (modern jen \textasciicircum{}) or *piar (modern pi t) (CKWP, ch. 8, pp. 6a, 1a); the graph ) or D might mean either *ngiwat (modern yüeh ♬) or *dziak (modern hsi ) (CKWP, ch. 7, pp. 6a, 7b-8a). Context would, in the first two cases at 67 ]}47 Similarly,
lated in the obligatory or (if a ch'i is present)
subjunctive mood. Thus, ``We should not build a
settlement (for if we do not) Ti will approve''
(拼编 147.2), matched by ``We might perhaps
build a settlement (but if we do) Ti will not assist
and approve (拼编 147.1); and ``On the next
ting-mao, we should not (perform a rain-seeking
dance, for if we do) it will not perhaps rain''
(拼编 442.7), matched by ``On the next ting-mao, we should perform a rain-seeking dance (for
if we do), it will rain” (拼编 442.6). I would
propose that two basic patterns are involved. The
first, more common one, is of the form: wu verb +
(for if we do verb) + negative result (e.g., 拼编 442.7, above; 拼编 12.2, translated in\ref{sec:chapters-ch03:3.7.1.2} (see 
sec. 3.7.1.2); 拼编 83.11; 96.26; 249.5; 300.2;
471.2; 易编 3787; 崔-编 1113 [the last two
S157.4]). The second, rarer pattern, is of the form:
wu verb (for if we do not verb) + positive
result (e.g., 拼编 147.2 above; 拼编
199.12). The general rule would be: when a wu
negative appears in the initial proposition, a
negative must appear either in the result expected
(the first pattern) or in the statement of hypothetical action supplied in parentheses (the second
pattern); i.e., if the result expected is negative,
the action supplied must be positive, and viceversa. And it follows that for the first pattern the
matching positive charge of the pair will be of
the form: verb + (for if we do verb) + positive
result; for the second pattern the matching positive
charge will be of the form: verb + (but if we do
verb) + negative result. These rules work for
all the cases cited above. For a different understanding of these patterns and their illocutionary
force, see Takashima (1976); (1977), pp. 11-18.
% source: scan 085, printed 67
the relative carelessness of the 商 engravers about details of graph shape, and the
inclusion or omission of structural features, denies us, in most cases, any opportunity to
connect epigraphic variations of the same graph with semantic or phonetic ones. Epigraphic
variation, when consistent, is generally related not to sound or sense, but to changes in
epigraphic style over time (table 26); far more rarely, it may be related to the positive and
negative phrasing of the charge, or to its right or left placement on the bone or shell.\footnote[48]{For examples, see the cases discussed in \ref{ch:5} (see ch. 5), n. 34. For one instance in which a graphic distinction appears to have had semantic, and perhaps phonetic, significance, see \ref{ch:1} (see ch. 1), n. 71.}48
Many 商 graphs can be analyzed in terms of five of the traditional six script
elements (liu shu): pictures (hsiang-hsing), abstract symbols (chih-shih), compound ideographs
(hui-yi), radical-plus-phonetic combinations (hsing-sheng or hsieh-sheng), and loans (chia-chieh) may all be identified. This suggests that the written script was already mature by
the historical period. Applied to the oracle-bone graphs, these classifications
are, of course,
frequently subjective and anachronistic. In many cases we lack the phonological information
to distinguish, say, a compound ideograph from a radical-plus-phonetic combination. And
some words were written with graphic variants of more than one script type.\footnote[49]{李 小庭 (1968), pp. 91-95, gives examples of each script type. Out of a sample of 1,226 graphs, he finds that 277 (23 percent) were pictures, 20 (2 percent) were abstract symbols, 396 (32 percent) were compound ideograms, 334 (27 percent) were radical-plus-phonetic, 129 (11 percent) were loans, and 70 (6 percent) could not be categorized. The fact that he found no instances of semantic extension (chuan-chu) may presumably be explained by our frequent inability to identify the original meanings of the 商 graphs that might have been ``extended.'' On the prevalence of radical-plus-phonetic graphs, cf. Barnard (1963), pp. 220-221.} 49 Further,
the loans are generally not true loans in which ``the character in question... does not have
its ordinary meaning but stands for another word, the reading and meaning of which it
adopts.\footnote[50]{Karlgren (1963), p. 2.} 50 Most of the so-called loan words can more properly be regarded as graphic
abbreviations. Thus IE is regarded as a loan for E, t as a loan for it, and as a loan
for\footnote[51]{These are among the 129 loans cited by 李 小庭 (1968), pp. 93-94. He subsequently argues (p. 101) that the bone graph was not a true loan for tsu, but a true picture, the radical being added later to avoid confusion with ch'ieh EL. But it is hard to be sure that many of his loan examples do not fall into the same category. For one case in which a picture graph (presumably its own phonetic) became in period V a radicalplus-phonetic with an additional signific-plusphonetic (?) added, see table 26, no. 15.} 51 Modern scholars must supply the radicals of later times in order to establish the
semantic value of the 商 word-such an approach is essential to deciphering the script
at all but it is anachronistic to suppose that the 商 wrote one word and meant
another. They simply wrote the word that way.
least, prevent any confusion, but such shorthand
notations suggest that we cannot always trust the
graphs for phonetic information. For the possibility that some graphs might have been notations
for a series of ritual actions, indicating, for example, the vessels and implements to be used, rather
than serving as word signs,[\textasciicircum{}52] consider the evidence
presented by Lefeuvre (1971); for a possible
Western Apache analogy, see Basso and Anderson
(1973).

[\textasciicircum{}52]: This approach cannot always yield certain results. That the sacrifice term (*mlộg, modern mao ) (\$507.4-508.3) was an abbreviation for *liôg (modern liu), "to slaughter, kill, cut up” (Serruys [1974], p. 92, n. 6), is phonetically and graphically plausible. But it remains only plausible until a full or transitional form of the graph is found in similar context.
% source: scan 086, printed 68
True loans which can assist scholars in their attempts to decipher the script are few.\footnote[53]{S577-588 lists some 320 cases of graphs which appear to have been used interchangeably. The majority, however, involve trivial epigraphic variations or abbreviations; some indicate semantic equivalence, if not identity (as with the disaster words at S579.1-580.1); some merely indicate topic substitution. Probably less than a dozen of the cases are phonetically useful. For a criticism of some of 岛 pairings, see 李 颜 (1969). 金 祥-heng (1967), p. 12a, also lists various ``loans'' in the inscriptions.} 53
The fact that *tieng □ (modern ting T ) and *tieng 7 (modern ting) were interchangeable
with the prefatory \# confirms the view that the verb ``to divine'' was pronounced like
*tieng and that the graph was a simplified depiction of a ting cauldron. Any attempt to
determine the meaning of the verb must take that phonetic and graphic evidence into
account.\footnote[54]{Cf. \ref{ch:2} (see ch. 2), n. 7. The relevant inscriptions are: \listitem{1} 钱-编 8.6.3 and 8.12.4; 易编 1834 and 钱-编 8.6.3 (the 钱-编 rubbing is hard to read; 岛 version is accepted by 李 颜 [1969], p. 197, but a different reading is given by Yeh Yü-sen [1934], ch. 8, p. 5a, and 李 Hsüeh-ch'in [1958], p. 53); and 易编 1010 (all S577.1); it may be noted that 李 学琴 (loc. cit.) speculates that all these bones were divined in the same seventh month; \listitem{2} 易编 8695 and 8888 (both from pit YH 251 and probably from different plastrons in the same set of five) and other examples at \$396.1-4. See too 拼编 381, k'ao-shih, pp. 445-446.}54
The equivalent usage or substitution of 3 (modern yu X) for + has also provided
valuable clues about the probable sound and meaning of 4, a 商 graph which has no
evident successor in 周 texts.\footnote[55]{For the relevant scholarship, see n. 85 and the note to table 26, no. 4. On the phonological problems see, in particular, Takashima (1973), pp. 30-31. The equivalence is particularly clear when the graphs were used as connectives in numbers, which might either be written +ELA “(ten years and five=) fifteen years'' (徐-编 1.44.5 [S197.1, but wrongly cited as 1.44.3]) or +2"(ten penned cattle and five=) fifteen penned cattle'' (崔-编 579). See the numerous examples in 陳 孟家 (1956), pp. 108I IO.}55 Similarly, the late-period substitution of 7 * (modern
wang tsai ) or ``, ``there will be no disaster,'' for the routine
tsai
(modern wang
) of the early inscriptions indicates a semantic and phonetic equivalence; all may
be read wang tsai, with the serving as phonetic in the late period graphs.\footnote[56]{For the relevant scholarship, see the note to table 26, no. 3.}56
The most famous case of what is thought to be phonetic substitution is one in which
the 商 engraver wrote *xwâr (modern huo \textasciicircum{}) instead of , ``disaster.'' Three
times on the same bone fragment the charge reads, ``in the next ten days there will be no
disaster,'' but in the one case it reads, “in the next ten days there will be no fire'' (fig.\footnote[57]{崔-编 1428 (\$579.1).} 19).57
If we assume that this was a scribal error rather than a case of deliberate topic substitution,
and that the scribal error may be explained in terms of phonetic identity (as well as possible
semantic equivalence), it may then be argued that must have been pronounced like
*xwar.\footnote[58]{Other evidence may give support. The phrase, ``Perhaps there will be fire'' (甲编 3083.1), appears analogous to the phrase ,"Perhaps there will be disaster'' (拼编 269.8, and other cases at S304.3-305.1). The of 易编 5349 (S303.2) may be analogous to the of 拼编 302.11-12 (\$173.3). Cf. Mickel (1976), pp. 72-75. 李 颜 (1967), pp. 21-26, also cites other cases of equivalence, but I cannot accept many of his readings. 69 0} 58 This accords fairly well with the usual interpretation of as *g'wa (modern
huo), ``disaster”. In later periods, on the other hand, , ``disaster,'' might also be
% source: scan 087, printed 69

written as a graph that meant ``bone,'' *kwǝt (modern ku) in period\footnote[59]{CKWT, p. 1497; S306.2-3.}\footnote[60]{CKWT, p. 3126; S217.4. On the interchangeability of and see 人文 1162, shakubun, p. 365, n. 1, and the inscriptions at \$217.4: 303.2.}\footnote[61]{天壤 k'ao-shih, pp. 7a-11a. As Mickel (1976), p. 71, notes, this hypothesis may be rejected on graphic as well as phonetic grounds.}\footnote[62]{Pulleyblank (1975), for example, depending on a revised table of 史-赤ng rhyme categories and the phonetic transformations undergone by words in related Sinitic languages, reconstructs the phonetic ``shape'' of the twenty-two cyclical graphs in Old Chinese and links them to the sounds and symbols of the Semitic alphabet. The oracle inscriptions themselves cannot be used to test such an approach. Pulleyblank argues, for instance, that the graphic element 7 represented the initial sound *s. The oracle-bone graphs *sien (or *sjien [周 Fa-kao (1973), 3484]) (modern hsin ) or *siang (modern shang ), which contain this element, presumably as a phonetic, support such a view. But reconstructions like *sien or *sjien were derived quite independently of any information in the oracle inscriptions.}\footnote[63]{Takashima \listitem{1973} provides the best introduction to the subject, the methodology to be used, and the difficulties involved. In this study he discerns semantic regularities in the initials, vowels, and finals of the negatives of the Wu Ting period (the phonetic reconstructions are his own): m *pjag(x)/pu不; *mjagx/wu 毋;\$*pjot fu 弗; and *mjat/wu . Following the lead of Serruys (1969; cf. 1974), he finds that *pjag(x) was used}
% source: scan 088, printed 70
How to Read the Graphs on a
Bone or Shell Fragment
Nothing appears more arcane, at first glance, than an oracle-bone rubbing. I
will try to demystify it in this section with a step-by-step analysis of how an actual turtle-shell piece, composed of two fragments (fig. 20), may be identified, deciphered, and placed
in context. The flow chart (fig. 21: summarizes the possible approaches.
How to Identify the Fragment
Orienting the piece is rarely a problem since rubbings are normally printed
correctly.\footnote[64]{Fragments have been printed upside down, especially in the earliest collections. See, for example, 铁-阴ün 23.4, 28.4, and the other cases listed in 颜 乙'ing's preface to the 乙-wen (Taipei, 1959) reprint of 铁-阴ün. Errors of this sort appear with some frequency in secondary works, e.g., Eisenberger (1938), p. 67 (the inscription is probably fake); 陳 知-mai (1966), fig. 4a (whole plastron inverted); William S. Y. Wang (1972), p. 53 (bottom two fragments, left column, inverted).} 64 If necessary, orientation is best done either by identifying certain graphs, or
graph elements, whose correct orientation is known, or by orienting the crack numbers
(near the top of the pu crack;\ref{sec:chapters-ch02:2.4} (see  sec. 2.4)) and crack notations (near the bottom of the pu
crack;\ref{sec:chapters-ch02:2.6} (see  sec. 2.6)) correctly. In figure 20, the crack number 6 and the shang-chi at the
left are readily placed. The pu cracks all run left; and the crack numbers 6, 7, and 8 run
right; these clues, together with the general configuration of the piece, lead us to suppose
that it comes from the right side of a shell in the region of the hyoplastron or hypoplastron
(fig.\footnote[65]{See \ref{ch:2} (see ch. 2), n. 122.} 3).65
As we have seen, inscriptions were written from top to bottom (sec. 2.9.4), but it is
necessary to determine in this case, as in all others, whether the inscription should be read
sideways and then down (in fig. 20, graphs 1, 4, 6, 2, 3, 5, 7; or 6, 4, 1, 2, 7, 5, 3) or read
down, with the columns read left to right (graphs 6, 7, 4, 5, 1, 2, 3) or right to left (graphs
I to 7 read in normal sequence).\footnote[66]{There is also the possibility that the graphs on a fragment do not belong to the same inscription (as defined at \ref{ch:2} (see ch. 2), n. 1) but are parts of two or more separate inscriptions; as we shall see, that is not the case here.}66 The fact that inscriptions generally run against the
cracks (sec. 2.9.4) suggests that the inscription in this case should be read from left to right,
but, as we see almost immediately from the position of the preface, we are dealing with a
common exception.\footnote[67]{See \ref{ch:2} (see ch. 2), n. 126. 0 71 U CO}67
The first, virtually automatic, step in identifying and translating any inscription is
basically for negating intransitive/passive verbs
and *pjat for negating transitive/causative verbs.
He also finds that the *P-type negative was used
when the situation involved natural forces or other
phenomena beyond 商 control, while the
*M-type was used when an element of human
will or control was involved. Nivison \listitem{1977} provides support for some of these conclusions. Mickel
\listitem{1976} has undertaken a similar study of the
``disaster'' words of the Wu Ting period (cf. \ref{ch:2} (see ch.
2), n. 80, above).
64. Fragments have been printed upside down,
especially in the earliest collections. See, for
example, 铁-阴ün 23.4, 28.4, and the other cases
listed in 颜 乙'ing's preface to the 乙-wen
(Taipei, 1959) reprint of 铁-阴ün. Errors of this
sort appear with some frequency in secondary
works, e.g., Eisenberger (1938), p. 67 (the inscription is probably fake); 陳 知-mai
(1966), fig. 4a (whole plastron inverted); William
S. Y. Wang (1972), p. 53 (bottom two fragments,
left column, inverted).
65. See \ref{ch:2} (see ch. 2), n. 122.
66. There is also the possibility that the graphs
on a fragment do not belong to the same inscription
(as defined at \ref{ch:2} (see ch. 2), n. 1) but are parts of two or
more separate inscriptions; as we shall see, that
is not the case here.
67. See \ref{ch:2} (see ch. 2), n. 126.
% source: scan 089, printed 71
to identify the cyclical graphs and, if there is one, the diviner's name recorded in the preface.
A glance at table 10 will indicate that graphs 1 and 2 may be read as chia-hsü. The
odds are that the next graph will be the diviner's name; graph 3 is, in fact, broken, but
we can see that it is composed of the roof element \textasciicircum{}, containing a horizontal line with
an inclined stroke below it. Studying the list of diviners' graphs at \$557-576, we find the
names of three diviners which contain a roof element: (S557), (S566), and (S572).
Of these, only the first seems to match the elements we have.\footnote[68]{The other two diviners may be eliminated by checking the single inscriptions (乙-chu 503 and 钱-编 6.31.1) in which they appear; neither of these is our fragment.} 68 The modern transcription
of this diviner's name, Pin or, may readily be obtained by consulting commentaries
(listed in table 9) to any of the inscriptions listed at\footnote[69]{E.g., Ikeda (1964), 1.2.15; 崔-编 36, k'ao-shih, p. 1ob; and 拼编, plate 29, k'ao-shih, p. 59, give a reading of . 践授 7.14, k'aoshih, p. 17b, and 甲编 59, k'ao-shih, p. 12, give賓.} \$557-558.69
So far, then, we know that the inscription starts: ``On chia-hsü (day 11), Pin. . .”
Since this part of the inscription lies near the edge of the fragment, it is probable that the
original preface was fuller, with a pu þ between graphs 2 and 3, just to the right of the
broken edge, and with a *tieng below graph 3 so that the restored inscription would
start: ``[裂] on chia-hsü (day 11), Pin [divined\footnote[70]{This is the prefatory formula most commonly used by Pin; see table 7, no. 1.1. Other possible formulas from the same table would be nos. 1.1.1, 1.4, 1.4.1, etc. Some of these, such as nos. 1.1.1 or 2.1, are unlikely, for they are RFG formulas. Pin was a period I court diviner. (His period may be learned from the glosses to 甲编, 拼编, etc., or from 焦 [1959], pp. 241-253; it may also be understood from the fact that the diviners at S557-576 are listed chronologically and that Pin is the first. On the court diviners and the RFG, \ref{sec:chapters-ch02:2.2.1.1} (see see sec. 2.2.1.1)).}]. . ``70
To identify the piece, it is necessary to select a key graph (that is, one that is not a
particle or frequently used, and for which 岛 is likely to provide a complete list of
occurrences) and locate it in Sōrui (sec. 3.3.2). In this case, graphs 5, 6, or 7 may all serve.
Starting with graph 5, we first separate its component elements. The bottom of the graph
is missing, but enough remains for us to identify the two elements as and M. 岛
radical chart (S590, or the endpapers) indicates that graphs containing either of these
elements are listed on page 4 of his index. Turning to page 4, it is more efficient to search
through the shorter list of graphs containing rather than the longer list containing \$.
The graph in question may be found in the third row from the bottom, the twelfth entry from
the right (the graph indexed is laterally inverted, **, but such rotation of graphs about
their central axis was common;\footnote[71]{Cf. \ref{ch:2} (see ch. 2), n. 125.}7¹ the word is the same despite the different graphic orientation). Thirteen inscriptions containing this graph are transcribed at S142.3 (reproduced
as fig. 22).
The most likely candidate is 易编 605, the only one to contain the right cyclical
date and the name of the diviner Pin. 岛 transcription supplies the pu, but this is not
justified by the rubbing; more seriously, his transcription does not contain characters 6
and 7, but, as the revised edition of the Sōrui \listitem{1971} indicates, 易编 605 has been joined
% source: scan 090, printed 72
with other fragments to form 拼编 390. A glance at that rubbing indicates that we
have indeed identified our two fragments-易编 603 + 605.72
How to Identify the Graphs
The revised edition of Sōrui \listitem{1971} contains modern transcriptions, together with
cross references to 甲-骨 wen-tzu chi-shih (sec. 3.3.1), for 840 of its entry heads. Thus, the
entry head for (S142.3) tells us that the graph is not found in 说文 and refers us to
CKWT, p. 3708, which, when we turn to that page, merely gives an incomplete list of
occurrences and the modern transcription chi with no discussion. Turning back to the
rubbings listed at S142.3 (fig. 22) we look for commentaries (table 9) published since
CKWT (1965). The 拼编 joinings of the 易编 fragments are the most promising lead.
On the theory that if 长 拼'üan glosses a graph he does so on its first 拼编
occurrence, we turn to 拼编 122, k'ao-shih, p. 182, where he identifies the graph as
chi Ẻ, the name of a place, probably to the west of the 商, used for hunting and farming.\footnote[73]{In fact, Ling-shih 113, k'ao-shih, p. 34b (published in 1959), identifies Chi as a place name. CKWT overlooked this gloss. Once we know its modern form, the identification of Chi as a place name may also be confirmed by consulting 焦 (1959), place name index, p. 27.} 73
That we are dealing with a place name is confirmed by graph 4. 岛 radical
chart indicates that is a radical in itself, the entries being given on S498 ff., where the
graph is identified as ts'ai (CKWT, p. 2049) or tsai \# (CKWT. p. 3991). Consulting
CKWT, p. 2049, we are told that ts'ai had the sense of tsai in oracle-bone writing.
Graphs 4 and 5, therefore, may be taken to mean “in Chi”.\footnote[74]{In the event that the 1971 edition of Sōrui, with its modern transcriptions and references to CKWT, is not available, one may arrive at the same conclusion by turning to a rubbing that has a modern transcription (table 9 lists such collections). Thus S498.4 might send us to 甲编 1132, which, as the k'ao-shih tells us, has been combined with 甲编 1110. The k'ao-shih for 1110 gives the modern form tsai \#. 11 J U 0 匚 ]}74
Graph 6, , is also one of 岛 radicals. The entries start on S289, where the
modern transcription t'ien is given. CKWT, pp. 4025-4026, quotes the 说文 defini-
72. The 拼编 tissue overlay labels the
fragments A, B, C, etc. The corresponding 易编
rubbing numbers are indicated at the bottom of
that overlay as well as at 拼编 399, k'ao-shih,
p. 454. In the event that only the first edition
\listitem{1969} of Sōrui, which does not indicate the
拼编 rejoinings, is available, further research
is needed to identify the fragment containing the
two missing graphs. Two routes are possible:
First, starting with graph 6, 岛 radical chart
indicates that is a radical itself and that the
entries start on S289. At S289.2 we find a list of
inscriptions containing the compound HI X, our
graphs 6 and 7. 易编 603 might be our missing
fragment. Since 易编 603 was assigned a rubbing
number close to that of 易编 605, a connection
is plausible. This possibility is supported by the
registration numbers 13.0.890 and 13.0.888
printed under the rubbings (for the meaning of
the registration numbers, see \ref{ch:1} (see ch. 1), n. 31), which
suggest that the fragments lay close together in
the pit, though on this point, cf. \ref{ch:4} (see ch. 4), nn. 184,
186. A study of the 易编 rubbings confirms that
the two fragments would fit together in the configuration we are trying to identify. Alternatively,
we may use the index to 易编 fragments rejoined
in 拼编 (Tang 阡-yüan [1969] and [1974])
to discover that 易编 603 and 605 were incorporated into 拼编 390.
% source: scan 091, printed 72
tion, to sow grain is called t'ien,'' and indicates that in the oracle inscriptions t'ien
generally had the meaning of ``hunt'' or ``field.''
Graph 7, \%, may be found under the “short-tailed bird'' radical on page 8 of the
Sōrui index, which leads us to S232.3 ff., where two modern transcriptions, huan and
kuan are given. 李 小庭 (CKWT, pp. 1299-1300) argues that both forms were
used for kuan, ``to observe."75
So far, then, we may transcribe the inscription as: 甲戍[][]」在娘」田觀.76
How to Decipher the Inscription
We may initially translate the charge as transcribed: ``In Chi, take to the field
and observe,'' with t'ien having the sense either of “hunt” or, following 说文, ``sow
crops.'' This translation must be refined by studying the context of related inscriptions.
Returning to S142.3 (fig. 22), we see that a majority of the inscriptions relating to Chi
contain the phrase. 岛 index (p. 10, row 10) sends us to S463.2, where the
entry head tells us that the first graph is shou or and where we are referred for this
compound to S194.2, which tells us that the second graph is nien ; over 250 occurrences
of the phrase are listed at S195.1 ff. Consulting CKWT, pp. 2366-2367, 2368, we find that
in the inscriptions nien has the sense of ``good harvest'' and that shou-nien is a prayer to
``receive harvest."77
The majority of the inscriptions which mention Chi at S142.3 (fig. 22) contain the
phrase ``at Chi receive harvest.'' None of them refers to the capture of animals. It is reasonable to assume, therefore, that Chi was a farming rather than hunting area and that t’ien
in the present inscription should be taken in its agricultural sense. The rightness of this
interpretation is confirmed by the headings at S232.4 where we find graph 7 forming
compounds with 1 (= shu, “millet, plant millet''), ¾ (=nien \#, ``harvest''), and t'ien .
Our tentative translation now reads, ``At Chi, sow crops and observe.''
But what does kuan (or huan?), graph 7, which we have translated as ``observe,''
mean? In questions of this sort, the right answer may require days of research. One possible
75. 李 also notes (CKIVT, p. 1301) that the
graph might be used as the name of a place or
person. That is the view adopted by 长 拼'üan in his transcription for 拼编 390.1,
k'ao-shih, p. 453, and he would apparently translate
as ``In Chi hunt (or cultivate [?]) Huan.'' But
there is no reason for thinking Huan to be in or
near Chi. The translation may be rejected.
76. Since we suspect that a * tieng is missing
below graph 3, it is also possible that one or more
graphs are also missing below graphs 5 and
7. In this case, however, 拼编 390, which
supplies the fragment below (易编 726), shows
there are no graphs missing from the charge.
77. Divinations about receiving harvest and
about other agricultural matters form one of the
categories (pu shou nien) of 焦's concordance (1959); see pp. 255-259 for the agricultural inscriptions of the diviner Pin. Our
inscription is not listed there because 焦 \listitem{1959}
was published before the appearance of 拼编
(1965). 焦, therefore, was working with the@@
two fragments 易编 603 and 605; one contained
the phrase t'ien kuan, with no identifying diviner
under which to classify it, the other contained the
diviner's name and tsai chi, but no classifiable
subject matter.
% source: scan 092, printed 73
solution might involve the following path. Turning first to the examples of interchangeable
graphs listed at the rear of Sōrui, we find (S586.2) * and * occurring in similar contexts
on two occasions. The fact that the double-mouth element might be omitted in oracle-bone
graphs had also been noted by Lo Chen-yü.\footnote[78]{Quoted by CKWT, p. 1297. Lo Chen-yü's opinion is also quoted by 甲编 444, k'ao-shih, p. 69, which could have been found by looking at 甲编 536.3, k'ao-shih, p. 82 (from S232.3).}\footnote[79]{Like Ikeda (1964), 人文, with its extensive glosses, is a major reference. Since a separate index volume gives the location of every graph appearing on the 3,256 fragments in the collection, it is not necessary to use Sōrui as a guide to its glosses. In this case, huan is one of the 540 说文 radicals by which the 人文 index is arranged. The graph may be found on p. x of the index volume, which directs us to p. 57 of the index text. The numbers in boldface type refer to 人文 entries where the graph is glossed; light type indicates its appearances without gloss.}\footnote[80]{The editors of CKWP had also proposed that referred to a kuan sacrifice. See their gloss to 崔-编 147 (CKWP, \ref{ch:4} (see ch. 4), p. 12b. col. 4).}\footnote[81]{One suggestion to make CKWT more manageable: toward the end of each entry. 李 小庭 gives his own conclusions about the various and often conflicting interpretations presented. These conclusions are signaled by a paragraph starting, in the standard manner of classical commentaries, with the character an ``I note...'' It is frequently best to start here and then work back through the entries which 李 discusses. KOV (XHDIH KODE Atto \#\textbackslash{}\textbackslash{}\textbackslash{}\textbackslash{}\textbackslash{} 俊上0元 X*1+簋orm中介A+5-0眾多 75 C ==}
% source: scan 093, printed 74
There is no sure Ariadne to guide us. CKWT, for instance, is not comprehensive;
人文's glosses are not complete and are out of date. No gloss is infallible. Scholars have
to plunge into the labyrinth for themselves, using the chart provided by figure 21. The
best guide despite occasional errors, variant transcriptions, and omissions-is 岛
Sōrui. Sōrui assembles the inscriptions; it guides us to CKWT and to the modern glosses;
it shows how a word was used in context. The rest we must do for ourselves. And we do it
by developing philological expertise, by studying the examples of graphic and semantic
evolution that commentators have proposed, and, above all, by reading as widely as
possible in the inscriptions themselves.

% source: scan 088, printed 70
\section[How to Read the Graphs on a Bone or Shell Fragment]{\origsecnum{3.6}How to Read the Graphs on a Bone or Shell Fragment}
\label{sec:chapters-ch03:3.6}

How to Read the Graphs on a
Bone or Shell Fragment
Nothing appears more arcane, at first glance, than an oracle-bone rubbing. I
will try to demystify it in this section with a step-by-step analysis of how an actual turtle-shell piece, composed of two fragments (fig. 20), may be identified, deciphered, and placed
in context. The flow chart (fig. 21: summarizes the possible approaches.

\subsection[How to Identify the Fragment]{\origsecnum{3.6.1}How to Identify the Fragment}
\label{sec:chapters-ch03:3.6.1}

Orienting the piece is rarely a problem since rubbings are normally printed
correctly.\footnote[64]{Fragments have been printed upside down, especially in the earliest collections. See, for example, 铁-阴ün 23.4, 28.4, and the other cases listed in 颜 乙'ing's preface to the 乙-wen (Taipei, 1959) reprint of 铁-阴ün. Errors of this sort appear with some frequency in secondary works, e.g., Eisenberger (1938), p. 67 (the inscription is probably fake); 陳 知-mai (1966), fig. 4a (whole plastron inverted); William S. Y. Wang (1972), p. 53 (bottom two fragments, left column, inverted).} 64 If necessary, orientation is best done either by identifying certain graphs, or
graph elements, whose correct orientation is known, or by orienting the crack numbers
(near the top of the pu crack;\ref{sec:chapters-ch02:2.4} (see  sec. 2.4)) and crack notations (near the bottom of the pu
crack;\ref{sec:chapters-ch02:2.6} (see  sec. 2.6)) correctly. In figure 20, the crack number 6 and the shang-chi at the
left are readily placed. The pu cracks all run left; and the crack numbers 6, 7, and 8 run
right; these clues, together with the general configuration of the piece, lead us to suppose
that it comes from the right side of a shell in the region of the hyoplastron or hypoplastron
(fig.\footnote[65]{See \ref{ch:2} (see ch. 2), n. 122.} 3).65

As we have seen, inscriptions were written from top to bottom (sec. 2.9.4), but it is
necessary to determine in this case, as in all others, whether the inscription should be read
sideways and then down (in fig. 20, graphs 1, 4, 6, 2, 3, 5, 7; or 6, 4, 1, 2, 7, 5, 3) or read
down, with the columns read left to right (graphs 6, 7, 4, 5, 1, 2, 3) or right to left (graphs
I to 7 read in normal sequence).\footnote[66]{There is also the possibility that the graphs on a fragment do not belong to the same inscription (as defined at \ref{ch:2} (see ch. 2), n. 1) but are parts of two or more separate inscriptions; as we shall see, that is not the case here.} 66 The fact that inscriptions generally run against the
cracks (sec. 2.9.4) suggests that the inscription in this case should be read from left to right,
but, as we see almost immediately from the position of the preface, we are dealing with a
common exception.\footnote[67]{See \ref{ch:2} (see ch. 2), n. 126.} 67

The first, virtually automatic, step in identifying and translating any inscription is
to identify the cyclical graphs and, if there is one, the diviner's name recorded in the preface.
A glance at table 10 will indicate that graphs 1 and 2 may be read as chia-hsü. The
odds are that the next graph will be the diviner's name; graph 3 is, in fact, broken, but
we can see that it is composed of the roof element ^, containing a horizontal line with
an inclined stroke below it. Studying the list of diviners' graphs at \$557-576, we find the
names of three diviners which contain a roof element: (S557), (S566), and (S572).
Of these, only the first seems to match the elements we have.\footnote[68]{The other two diviners may be eliminated by checking the single inscriptions (乙-chu 503 and 钱-编 6.31.1) in which they appear; neither of these is our fragment.} 68 The modern transcription
of this diviner's name, Pin or, may readily be obtained by consulting commentaries
(listed in table 9) to any of the inscriptions listed at\footnote[69]{E.g., Ikeda (1964), 1.2.15; 崔-编 36, k'ao-shih, p. 1ob; and 拼编, plate 29, k'ao-shih, p. 59, give a reading of . 践授 7.14, k'ao-shih, p. 17b, and 甲编 59, k'ao-shih, p. 12, give賓.} \$557-558.69

So far, then, we know that the inscription starts: ``On chia-hsü (day 11), Pin. ...''
Since this part of the inscription lies near the edge of the fragment, it is probable that the
original preface was fuller, with a pu þ between graphs 2 and 3, just to the right of the
broken edge, and with a *tieng below graph 3 so that the restored inscription would
start: ``[裂] on chia-hsü (day 11), Pin [divined\footnote[70]{This is the prefatory formula most commonly used by Pin; see table 7, no. 1.1. Other possible formulas from the same table would be nos. 1.1.1, 1.4, 1.4.1, etc. Some of these, such as nos. 1.1.1 or 2.1, are unlikely, for they are RFG formulas. Pin was a period I court diviner. (His period may be learned from the glosses to 甲编, 拼编, etc., or from 焦 [1959], pp. 241-253; it may also be understood from the fact that the diviners at S557-576 are listed chronologically and that Pin is the first. On the court diviners and the RFG, \ref{sec:chapters-ch02:2.2.1.1} (see see sec. 2.2.1.1)).}} 70]. . ``

\subsection[How to Identify the Graphs]{\origsecnum{3.6.2}How to Identify the Graphs}
\label{sec:chapters-ch03:3.6.2}

To identify the piece, it is necessary to select a key graph (that is, one that is not a
particle or frequently used, and for which 岛 is likely to provide a complete list of
occurrences) and locate it in Sōrui (sec. 3.3.2). In this case, graphs 5, 6, or 7 may all serve.
Starting with graph 5, we first separate its component elements. The bottom of the graph
is missing, but enough remains for us to identify the two elements as and M. 岛
radical chart (S590, or the endpapers) indicates that graphs containing either of these
elements are listed on page 4 of his index. Turning to page 4, it is more efficient to search
through the shorter list of graphs containing rather than the longer list containing \$.
The graph in question may be found in the third row from the bottom, the twelfth entry from
the right (the graph indexed is laterally inverted, **, but such rotation of graphs about
their central axis was common;\footnote[71]{Cf. \ref{ch:2} (see ch. 2), n. 125.} 71 the word is the same despite the different graphic orientation). Thirteen inscriptions containing this graph are transcribed at S142.3 (reproduced
as fig. 22).

The most likely candidate is 易编 605, the only one to contain the right cyclical
date and the name of the diviner Pin. 岛 transcription supplies the pu, but this is not
justified by the rubbing; more seriously, his transcription does not contain characters 6
and 7, but, as the revised edition of the Sōrui \listitem{1971} indicates, 易编 605 has been joined
with other fragments to form 拼编 390. A glance at that rubbing indicates that we
have indeed identified our two fragments-易编 603 + 605.72

How to Decipher the Inscription

We may initially translate the charge as transcribed: ``In Chi, take to the field
and observe,'' with t'ien having the sense either of "hunt" or, following 说文, ``sow
crops.'' This translation must be refined by studying the context of related inscriptions.
Returning to S142.3 (fig. 22), we see that a majority of the inscriptions relating to Chi
contain the phrase. 岛 index (p. 10, row 10) sends us to S463.2, where the
entry head tells us that the first graph is shou or and where we are referred for this
compound to S194.2, which tells us that the second graph is nien ; over 250 occurrences
of the phrase are listed at S195.1 ff. Consulting CKWT, pp. 2366-2367, 2368, we find that
in the inscriptions nien has the sense of ``good harvest'' and that shou-nien is a prayer to
``receive harvest.''77

The majority of the inscriptions which mention Chi at S142.3 (fig. 22) contain the
phrase ``at Chi receive harvest.'' None of them refers to the capture of animals. It is reasonable to assume, therefore, that Chi was a farming rather than hunting area and that t'ien
in the present inscription should be taken in its agricultural sense. The rightness of this
interpretation is confirmed by the headings at S232.4 where we find graph 7 forming
compounds with 1 (= shu, "millet, plant millet''), ¾ (=nien #, ``harvest''), and t'ien .
Our tentative translation now reads, ``At Chi, sow crops and observe.''

But what does kuan (or huan?), graph 7, which we have translated as ``observe,''
mean? In questions of this sort, the right answer may require days of research. One possible
solution might involve the following path. Turning first to the examples of interchangeable
graphs listed at the rear of Sōrui, we find (S586.2) * and * occurring in similar contexts
on two occasions. The fact that the double-mouth element might be omitted in oracle-bone
graphs had also been noted by Lo Chen-yü.\footnote[78]{Quoted by CKWT, p. 1297. Lo Chen-yü's opinion is also quoted by 甲编 444, k'ao-shih, p. 69, which could have been found by looking at 甲编 536.3, k'ao-shih, p. 82 (from S232.3).}} 78

There is no sure Ariadne to guide us. CKWT, for instance, is not comprehensive;
人文's glosses are not complete and are out of date. No gloss is infallible. Scholars have
to plunge into the labyrinth for themselves, using the chart provided by figure 21. The
best guide despite occasional errors, variant transcriptions, and omissions-is 岛
Sōrui. Sōrui assembles the inscriptions; it guides us to CKWT and to the modern glosses;
it shows how a word was used in context. The rest we must do for ourselves. And we do it
by developing philological expertise, by studying the examples of graphic and semantic
evolution that commentators have proposed, and, above all, by reading as widely as
possible in the inscriptions themselves.

\section[A Plastron Set]{\origsecnum{3.7}A Plastron Set}
\label{sec:chapters-ch03:3.7}

A Plastron Set
Translation not only involves placing inscriptions in the context of other inscriptions. It also involves understanding the role that divination and the carving of the inscriptions played in 商 culture; it involves understanding, as far as we can, what the diviner
thought he was doing.\footnote[82]{For a preliminary analysis of these issues, which will be considered more fully in Studies, see Keightley (1973a) and (1975c). Cf. too, Takashima (1976), who proposes a rationale for some of the charges in the set discussed below.} 82 It is instructive in this regard to examine one set of inscriptions
in detail in an initial attempt to explore the rationale underlying 商 divination and to
indicate the kinds of questions that we need to ask.
I select the set of inscriptions carved on the five plastrons 拼编 12-21 because,
when the charges on the front are considered together with the subcharges (sec. 3.7.4.2)
and prognostications on the back, they form one of the most informative multiple-plastron
sets we now possess. (The fact that these five plastrons form a set [sec. 2.5] is demonstrated
by the similarity of the divination charges and the ordinal sequence of crack numbers
[sec. 2.4], which are all 1 on 拼编 12, all 2 on 拼编 14, and so on.)
The analysis which follows should, if possible, be studied with the 拼编 rubbings
and 长 并-chüan's k'ao-shih at one's side; I do not always follow his interpretation,
and I make no attempt to reproduce his extensive notes about the identity of the various
figures and statelets involved. My concerns are primarily with the divination process.

\subsection[The Inscriptions]{\origsecnum{3.7.1}The Inscriptions}
\label{sec:chapters-ch03:3.7.1}

\subsubsection[Duplication and Abbreviation]{\origsecnum{3.7.1.1}Duplication and Abbreviation}
\label{sec:chapters-ch03:3.7.1.1}

With the exception of minor variations in layout and epigraphy, inscriptions I,
2, 5, 8, 9, and 10 on the front of the shells and 1 through 8 on the back were ``carbon
copies,'' reproduced without abbreviation on all five plastrons.\footnote[83]{This conclusion depends, in several cases, on supplying graphs that would, it is assumed, have appeared on missing shell fragments. Due to the roughness of the inner surface of the shell, it is not always easy to see the graphs in the rubbings of the plastron backs; in these cases I rely on the k'ao-shih.}83 The other inscriptions
% source: scan 094, printed 75
on the front abbreviate the preface or the name of 知 郭 but with no evident regularity.
The general impression is that there was a desire to reproduce each inscription fully on
each plastron but that it was not essential to do so.\footnote[84]{Dealing with the fronts of the shells and supplying missing graphs in each case: sentence 3 was recorded with nine graphs on 拼编 12 and 14, with five graphs on 拼编 16 and 20 (the kan-chih date, diviner's name, and pu graph being omitted), and with eight graphs on 拼编 18 (知 郭's name being abbreviated to 知); sentence 4 was recorded with ten graphs on 拼编 14, the second plastron in the set, but with six graphs on the other four shells (which omit the kan-chih date, the diviner's name, and pu graph; the k'ao-shih for 拼编 18.4 is in error); sentence 6 was recorded with eleven graphs on all but 拼编 12, where 知 郭 was again abbreviated to 知; sentence 8 was recorded with seven graphs on all but 拼编 12, where a one-graph preface, *tieng, was added to make a total of eight graphs. Ranked in order of completeness, 拼编 14, the second plastron in the set, has ninety-eight graphs; 拼编 12, the first plastron, has ninety-four; 拼编 18, the fourth plastron, has ninety-three; and 拼编 16 (fig. 15) and 20, the third and fifth plastrons respectively, both record the same ninety graphs. But even though 拼编 14 has the fullest record, 拼编 12 has a fuller record (by one graph) of sentence 8. No principle appears to explain this pattern of abbreviation; was apparently due to carelessness. On the various kinds of abbreviation employed by the 商 engravers, see 长 D 77 44 ] D C C C C C C C C C C C C C D CO}84
The records about the source of the shells or the results of the divination were not
duplicated or abbreviated in the same way as the charges. Thus, the marginal notations
(拼编 13.9; 21.9) appear only on the back of the first and last shells; and the two
prognostications appear only on 拼编 17 and 19, the back of the third and fourth shells.

\subsubsection[Translation]{\origsecnum{3.7.1.2}Translation}
\label{sec:chapters-ch03:3.7.1.2}

Given the identity or similarity of the divination sentences, we may provide a
combined translation for the set as a whole. Taking the inscriptions on 拼编 16/17
(figs. 15 and 16), the most complete plastron, to be representative, and supplying graphs
from the other four plastrons as necessary, I transcribe and translate the full text of the
inscription set as follows:
Reconstructed Chinese Text
for 拼编 12, 14, 16. 18, 20
(I) 辛酉卜殻\#今王从望乘伐下危受ㄓㄨ」\listitem{2} 辛酉卜殻\#今王勿从望乘伐下危弗其受ㄓㄨㄥ
\listitem{3} 辛酉卜段件(今)王从沚(伐巴方受ㄓㄨ)儿\listitem{4} 辛酉卜段(今)王勿从(伐巴方弗其受ㄓㄨ)」\listitem{5} 辛酉卜段网(今)王\listitem{6} 辛酉卜殼(今)王勿从(伐巴方受ㄓ文)
让从(伐巴方弗其受ㄓㄨ)」\listitem{7} ㄓ犬于父庚卯羊L
\listitem{8} 祝忄之疾齒鼎龍\listitem{9} 疾齒龍(10)(疾齒)不其龍
% source: scan 095, printed 76
Translation
of 拼编 12, 14, 16, 18, 20
(Preface:)
(Charge:)
(Preface:)
(Charge:)
PAIRS OF COMPLEMENTARY CHARGES
Negative charge (left side)
\listitem{2} 裂 on hsin-yu (day
58), Ch üeh divined:
``This season, the king should not
follow Wang 成 to attack the 夏
韦. for if he does, we) will not perhaps receive assistance in this case.\footnote[85]{长 拼'üan reads the bone graph as ch'un \#, ``spring.'' 周 Feng (1976), p. 222, reverting to an earlier theory of Yang Shu-ta , argues that the graph (composed, in his view, of k'ou □ plus ts'ai *) was a loan for tsai ``year.'' CKWT (pp. 1972-1974) rejects both readings and argues, with Yü 邢-wu (1940), pp. 6a-8b, that the graph was ancestral to *d'iôg or *t'iôg t'iao, which is close in sound to *ts'iôg/ch'iu K, ``autumn.'' A study of all the inscriptions containing S187.2-4) reveals that the period to which it applied extended from at least the eleventh to the fifth 商 month. ``Season'' is offered as a functional translation. For the reason why I translate these charges in the obligatory mood, see n. 44. The phrase shou yu yu has been much discussed. There is little doubt that the final graph stands for yu, thiori, ``(divine) assistance.'' As to the 4 (or, in later oracle texts, * ; see table 26, no. 4.3), I follow the view that it has an attributive, pronominal function. This is the solution briefly suggested by scholars such as Wu 奇-ch'ang ([1971], pp. 10It) and Ch’en 孟家 ([1956], pp. 96-97); it is well and extensively supported by David Nivison (1977). In this view, the shou yu yu of this inscription would mean ``receive this assistance'' (i.e., assistance in this attack against the 夏 韦); shou yu nien would mean ``receive this harvest'' (i.e., the one divined about, as in fig. 9). For the alternative theory (discussed by Nivison [1977], pp. 18-19) that 4 in such phrases stands for yu \#, ``abundant,'' see Takashima (1973), pp. 336-338; cf. Serruys (1974), pp. 29-30. For additional scholarship, see 金 祥-heng (1967); Nivison (1973)\%; Yen Yi-p'ing (1973a).}\footnote[86]{I assume, on the basis of context, that sentences 3 and 4 are in the same conditional mood as I and 2, though recorded in abbreviated form. This is speculative. If the meaning in parentheses is not supplied, then the sentences should read: ``The king should follow 知 Kuo/The king should not follow 知 Kuo.'' For ``should not'' as a translation for the negative wu, see n. 44. To supply the Pa-fang as the object of the attack is still more speculative. I do so for the following reasons. No inscription (S65.4 ff.) records that 知 Kuo was to attack the 夏 韦. It is unlikely, therefore, that sentences 3 to 6 should be taken as proposing that the king join 知 Kuo in attacking the 夏 韦. 并-编 22.11-12, probably divined by Ch'üch on keng-shen (day 57), concerns the king following 知 Kuo to attack the Pa-[fang]. 并-编 24 + 京-chin 1266 (see 并-编 24, k'ao-shih, pp. 48-51) records eight charges (A-H), probably all divined by 鄭 on hsin-yu (day 58), about the king following Wang 成 to attack the 夏 韦 and twelve charges (I-N, P-U),}"85
\listitem{4} 裂 on hsin-yu (day
58), Ch'üeh divined:
``(This season) the king should not
follow 知 郭 (to attack the 八. for if he does, we will not perhaps
receive assistance in this case)."86
Positive charge right side)
\listitem{1} 裂 on hsin-yu (day
58), Ch'üch divined:
``This season, the king should follow
Wang 成 to attack the 夏 韦,
(for if he does, we will receive assistance in this case.''
\listitem{3} 裂 on hsin-yu (day
58), Ch'üeh divined:
``(This season) the king should follow
知 郭 (to attack the 八, for
if he does, we will receive assistance in
this case).''
拼'uan (1956), pp. 246, 253-254; 周
洪祥 (1969), pp. 60-69; 李 大
(1972), pp. 134-169. 李 notes (pp. 137, 162-163)
that it is the later inscriptions in a set that are
likely to be abbreviated; but that is not true in
this case. He also notes pp. 138, 157) that when
three or more inscriptions refer to the same topic,
the preface and charge are both frequently
abbreviated. It is probable that in extreme cases,
the entire inscription could be abbreviated out of
existence; see \ref{ch:2} (see ch. 2), n. 32, above, on inscriptionless
cracks.
% source: scan 096, printed 77
(Preface:)
(Charge:)
(Preface:)
(Charge:)
\listitem{6} 裂 on hsin-yu (day
58), Ch'üch divined:
``(This season) it should not be 知
郭 that the king follows (to attack the
八, for if he does we will not perhaps receive assistance in this case).\footnote[87]{On the copulas and, see Takashima (1973), pp. 246–253, 399, n. 1; Serruys (1974), pp. 24-25, 74, n. 64; see too, Takashima (1977), pp. 11-18, for possible contrasts of emphasis involved in the different sentence patterns. If the meaning in parentheses is not supplied, then sentences 5 and 6 should be translated: ``It should be 知 郭 whom the king follows / It should not be 知 郭 whom the king follows.'}\footnote[88]{On the mao sacrifice, see n. 52. ``}\footnote[89]{The translation of sentence 8 is problematical. On the verb of, see \ref{ch:1} (see ch. 1), n. 62. The graph is graphically distinct from the M normally used in the preface of the period I 79 U B SEC C C D U ]}
% source: scan 097, printed 78
(Subcharge:) \listitem{1}
(Subcharge:) \listitem{3}
(Subcharge:) \listitem{5}
(Subcharge:) \listitem{7}
A PLASTRON SET
Translation
of 拼编 13, 15, 17, 19, 21
PAIRS OF COMPLEMENTARY CHARGES
Negative charge (left side)
It is because of Fu 甲.
It is because of Fu Keng.
It is because of Fu 新.
It is because of Fu 乙.
Positive charge (right side)
\listitem{2} It is not because of Fu 甲.\footnote[90]{For wei in the sense of ``due to,'' ``because of,'' ``caused by,'' see Takashima (1973), pp. 243, 394,n.34.}90
\listitem{4} It is not because of Fu Keng.
\listitem{6} It is not because of Fu 新.
\listitem{8} It is not because of Fu 乙.

(First prognostication,
17.9:*
(Second prognostication,
19.9:
(Marginal notations,
13.9; 21.9:)
NONCOMPLEMENTARY RECORDS
The king [reading the cracks] said: ``There will perhaps be a natural
phenomenon. If it be a wu day when there is a natural phenomenon,
it will be inauspicious."92
The king, reading the cracks, said: ``On ting-chou (day 14), there will
perhaps be a
natural phenomenon; it will be inauspicious. If it be a
chia day when there is a natural phenomenon, it will be auspicious.
If it be a hsin day when there is a natural phenomenon, it will also
be inauspicious.93
知 Fei (?) sent in two (shells). In Tuan (?).94
* Note that my numbering and reconstruction of the prognostications differ from 长 拼'üan's;
see n.\footnote[91]{I believe 长 拼'üan errs in separating his 拼编 17.9 and 17.10, in putting his 17.9 before his 17.10, and in the reconstruction of his 17.10 (拼编 17, k'ao-shih, p. 39). For parallel inscriptions to justify my arrangement and translation, see 钱-编 7.32.3; 徐-编 5.9.2, etc. (S94.2). A modern equivalent of has not been identified. For a resumé of the scholarship, see 拼编 17, k'ao-shih, p. 40; 人文 3, shakubun, p. 143, n. 2; Serruys (1974), p. 76; Mickel (1976), pp. 149-158. In some contexts the graph refers to a sacrifice (e.g., Kikkō 1.26.6 [S94.2]; 追 286 [S94.3]). But in other instances it refers to some natural, probably climatological, event whose significance was thought to vary with the day of its occurrence.} 91.
inscriptions (see \ref{ch:2} (see ch. 2), n. 7). For other inscriptions
in which this ornate form of ting was used as a
verb, apparently of sacrifice, see 京 3875;
4330; 侯编 1.6.4; 乙-村'un 783 (all \$396.2).
Ting cauldrons were used for offering ancestral
sacrifices in 周 times. In the oracle bones the
graph may simply have had the generic sense of
``the ritual performed at the cauldron.'' CKWT,
pp. 3477-3484, takes the view that 5 and were
the same word; they appear in different contexts,
however, and 岛 properly provides separate
entries (S241.1 ff.; 242.1 ff.). Scholars are generally
agreed that our graph, 5, is an early form of lung
长 拼'üan (拼编 12, k'ao-shih, p.
32) takes it as loan for hsiung X, but this may be
rejected for two reasons: First, following Karlgren's studies of initials (1963), p. 10, it is unlikely
that liung and *xiung would have been
interchangeable in this way. Second, in the
inscriptions at S242.1 ff., the phrase is
common; following the rule of undesirability (see
n. 41), which views the ch'i as a particle that
weakens the undesirable charge, the lung of ``not
perhaps lung'' would be a desirable rather than
undesirable word. Lung was probably used as a
loan for *t'liung/ch'ung, ``grace, favor,'' as it was
used in 史-赤ng, ode no. 293 (GSR, 1193), and
so I take it here; cf. 焦 (1959), p. 138; 郭
沫若, cited by CKWT, p. 3477. Nivison's argument ([1977], p. 12) that 5, which following
岛 ([1958], p. 278) he identifies as ch'iu
``hornless dragon,'' was ancestral to ch'ou,
``recover,'' should also be considered.
% source: scan 098, printed 79

\subsection[The Hollows and Cracks]{\origsecnum{3.7.2}The Hollows and Cracks}
\label{sec:chapters-ch03:3.7.2}

Supplementing our ignorance about the piece missing from the left epiplastron
of 拼编 17 with information derived from the corresponding areas of 拼编 15, 19,
and 21, we can infer that each plastron had thirty-two hollows chiseled into its back, with
sixteen on the right and sixteen on the left, arranged in a symmetrical pattern (fig. 23).
A total of 160 hollows was made, therefore, on the five plastrons. A study of the pu cracks
appearing on the front of the plastrons suggests that the hollows were all burned and
cracked; the hollows were used with full efficiency (cf. sec.\footnote[92]{For ``if'' as a translation of ch'i in prognostications, see n. 41.}\footnote[93]{The hsin is illegible in the rubbing. Here I follow the transcriptions of 长 拼'uan and S94.2.}\footnote[94]{Following 拼编 13, k'ao-shih, p. 36, 知 (or Sui 7 [?]-see CKWT, pp. 18951896) is a name that appears in four uninformative period I inscriptions (S61.2). Fei (?) is a hapax legomenon. Possibly, 知 Fei referred to a leader called Fei from the region of 知. I assume with 岛 (S357.4) that is a variant of . 长 拼'üan, in fact, transcribes the graph as ; this is possible on the basis of 拼编 13, but 岛 reading is supported by 拼编 21. The roughness of the plastron back makes the accuracy of the rubbings uncertain. As a matter of convenience, I follow Ikeda (1964), 2.15.14 and 2.32.14, for the highly speculative reading of Tuan (?). On this place, see n. 98 below.}\footnote[95]{The rubbings do not allow certainty on this point for they do not always reproduce the pu cracks, especially when the cracks were not engraved (cf. 长 拼'üan [1956], p. 241). Note, for example, that 拼编 20.3 has a crack number 5, but no crack is visible in the rubbing; similarly, the hsiao-chi crack notation on the left hypoplastron of 拼编 18 lacks a visible crack. The fact that cracks can be identified in each of the thirty-two hollow positions for at least one of the plastrons suggests that all the hollows were cracked in each case. This analysis, and much of what follows, depends on the use of tracings made of the hollow positions, reversed and superimposed on the front of the shell, to locate the area where cracks should have formed and near which the relevant inscriptions should have been carved. more}\footnote[96]{For example, two cracks without numbers appear at the top center of the right and left hypoplastrons of 拼编 16, and two appear at the top of the right and left xiphiplastrons; the cracks associated with 拼编 18.7 and 18.8 have no numbers (the k'ao-shih appears to be wrong), and at least two more inscriptionless cracks on the right hyoplastron of 拼编 18 lack numbers; the cracks for 拼编 20.7 and 20.8 lack numbers (cf. the k'ao-shih, p. 43), and so do at least five more inscriptionless cracks on the right and left hyo- and hypoplastrons of that shell. It is possible that in some of these cases the numbers may be on fragments now lost, but it is clear that numbers were not carved in every case (cf. \ref{ch:2} (see ch. 2), n. 32, above). It appears that at the top center of the left hypoplastron of 拼编 14 the 商 engraver deliberately obliterated a crack number (k'ao-shih, p. 37); on this practice, \ref{sec:chapters-ch02:2.4} (see see sec. 2.4). CCC C D U 0 81 \}}
% source: scan 099, printed 80
Lacking any direct evidence, we must assume that the rule of proximity determined
which inscriptions were related to which cracks. As we have seen, inscriptions appear to
have been carved above, or to the side of, their pu cracks and on the side of the crack which
lacked the transverse branch (sec. 2.9.4). Using this criterion, we can separate the thirty-two cracks formed in each plastron into three categories: \listitem{1} front-inscription cracks,
which may be linked by their location to one of the charges carved on the plastron front;
\listitem{2} back-inscription cracks, which may be linked by their location to one of the charges
carved on the plastron back;\footnote[97]{The examples cited by 李 大 (1972), pp. 69-79, clearly indicate that charges on the back of a shell were related to cracks that appeared on the front.} 97 and (3 inscriptionless cracks, which can be linked by
their location to no adjacent inscriptions. The hypothetical relation between the cracks
and inscriptions is given in table 11. The fact that in this view Fu Keng is associated with
only two cracks is anomalous; the other ancestors are all associated with four. And we
cannot be sure that the eight remaining hollows at the bottom of the plastron were not
also cracked for Fu 乙. But I believe that as a rough approximation these correlations are
correct. Of the thirty-two cracks made on each plastron, ten were front-inscription cracks
and fourteen were back-inscription cracks, leaving eight inscriptionless cracks at the
bottom of the shell.
Still more speculatively, I assume that the prognostications at the top of 拼编 17
and 19-breaking the rule of proximity-are to be correlated with those eight seemingly
inscriptionless cracks. This could be justified by a rule of balance in which the top and
bottom of the shell stood in special relation to each other, just as the sides did. And the
distance between the hollows and their related inscriptions could be explained functionally:
the necessary cracks could not have been made on the epiplastron or entoplastron at the
top of the shell because they would have broken into the zone of inscriptions 17.1, 17.2, 19.1,
and 19.2 on the front; and the prognostications could not have been carved on the xiphiplastron at the bottom of the shell because hollows H25 to H32 (see fig. 23) left too little
room for such lengthy inscriptions. I tentatively assume, therefore, that these last eight
cracks were associated not with recorded charges but with eight unrecorded subcharges
whose existence we may also infer from the resulting prognostications (\ref{sec:chapters-ch03:3.7.4.2} (see see sec. 3.7.4.2)).

\subsection[The Crack Notations]{\origsecnum{3.7.3}The Crack Notations}
\label{sec:chapters-ch03:3.7.3}

There are three crack notations on the front of the plastron pieces now remaining
to us. The first is a hsiao-chi, ``slightly auspicious,'' carved in the lower angle of the
front-inscription crack formed in hollow H6 of 拼编 14/15; I have associated this (in
table 11) with charge 8: ``Praying to lead away this sick tooth, the ting-sacrifice will be
favorable.'' The second crack notation is a shang-chi, ``highly auspicious,'' near the
top left corner of the right hypoplastron of 拼编 14; this notation lies in the lower
% source: scan 100, printed 81
angle of the crack formed in hollow H15, which I have associated with subcharge 5: ``It
is because of Fu 新.'' The third crack notation is another hsiao-chi, ``slightly auspicious,''
just visible on the fragmented edge of the left hypoplastron of 拼编 18; this notation
lies in the lower angle of the crack in hollow H16, which I have associated with subcharge
6: It is not because of Fu 新.''
C C C D D

\subsection[Interpretation]{\origsecnum{3.7.4}Interpretation}
\label{sec:chapters-ch03:3.7.4}

Having tentatively correlated the cracks, crack notations, and inscriptions, we
may now attempt to reconstruct the divination scenario. Two plastrons, the first and last
of the set, were sent in by 知 Fei at a place called Tuan; this fact was recorded as a
marginal notation. Other inscriptions reveal that the 商 farmed and sojourned at
Tuan and that it was the site of a garrison camp.\footnote[98]{See 拼编 169.5-6; 易编 1044; 1410 (S357.4). Since the two shells of our set were received at Tuan, it is conceivable that the divinations they record about the season's campaigns against the 夏 韦 and the 八 were divined at ``camp Tuan,'' to which the king may have gone in the second month (拼编 65.8 [S357.4]). The lack of any notation to the effect that our plastrons were ``ritually prepared'' might be explained by the fact that none of the royal women who usually performed this task (sec. 1.4) had accompanied the king to Tuan, where the shells were received.} .98 As a locus of such activities it was natural
that certain shipments of shell (sec. 1.2.4) would have been received there.\footnote[99]{For other marginal notations recording the receipt of small numbers of shells at various places, see S256.}99 The source
of the other three shells in the set, one of which (拼编 16/17) has been identified as
Ocadia sinensis (tables 3 and 4), was not recorded.\footnote[100]{Another group, of three shells, was apparently received at Tuan (易编 642 [S357.4]), but the fragment recording that fact does not appear to belong to our set.}100
After the marginal notations had been carved in the shell (sec. 1.4), thirty-two double
hollows (sec. 1.5) were chiseled in the back of each of the five plastrons, which were then
ready for cracking. If we assume that it required thirty minutes to chisel a double hollow,
eighty man-hours would have been required to prepare the hollows on the five plastrons.\footnote[101]{See the report of 长 Kuang-yüan, cited in \ref{ch:1} (see ch. 1), n. 93.} 101

\subsubsection[The Front of the Plastron]{\origsecnum{3.7.4.1}The Front of the Plastron}
\label{sec:chapters-ch03:3.7.4.1}

The first six charges concern contemplated attacks against two enemy statelets, the
夏 韦 and the 八.\footnote[102]{In the case of the 夏 韦 F, the reading of the second 商 graph, \&, is uncertain; for an introduction to the scholarship and the location of the statelet, see Yü 邢-wu (1940), pp. 22a-b; 岛 (1958), pp. 389-390; Ikeda (1964), 1.18.3; Serruys (1974), p. 107, n. 39. In the case of the 八, there are disagreements about the reading of the first 商 graph, , and the location of the statelet; for an introduction to the problems, see 陳 孟家 (1956), p. 284; 岛 (1958), pp. 390-391; CKWT, p. 2783.} 102 The question in the king's mind was whether he should
follow Wang 成\footnote[103]{Wang 成 was a major 商 general or ally in period I. For an introduction to the scholarship, see Ogawa 1, shakubun, pp. 250-251. 83 CC C 0 U U}103 in his attack on the 夏 韦 or whether he should follow 知
% source: scan 101, printed 82
郭 104 in his attack on the 八. (It is important to stress that the words “to attack
the 八'' are recorded nowhere in this set. If I am correct in supplying this meaning,\footnote[104]{知 was a place name. 郭 (?)—hence 郭 of 知? is thought to have been a powerful general of period I (e.g., CKWT, p. 3545). The reading of kuo for is not certain. Some scholars (e.g., 长 拼'üan, 拼编 12, k'ao-shih, p. 30) read it as chia . For an introduction to the scholarship, see 岛 (1958), pp. 387, 391; Ikeda (1964), 1.17.5; Serruys (1974), pp. 113-114.}\footnote[105]{For the reasons why I do so, see n. 86.}\footnote[106]{For the role of Ti in such divinations, see the inscriptions at S157.4, especially 易编 3787. Cf. 胡侯-宣üan (1959), pp. 37-39.}\footnote[107]{There are reasons for thinking that the positive charge was cracked before the negative (see \ref{ch:2} (see ch. 2), n. 130), but it is possible that the negative charge was cracked first, that the two charges were cracked simultaneously, or that there was no consistent rule.}\footnote[108]{At least ten cracks must have been made on keng-shen (day 57) for 徐-编 3.11.3; 3.12.1; 徐-村'un 1.624; 崔-编 1108; 1109 (all S110.2). Twenty more cracks were formed for the set represented by 拼编 22.11-14. On hsin-yu (day 58), in addition to the thirty cracks of our set, some one hundred cracks must have been made for the set represented by 拼编 24 + 京 1266. It is not certain, of course, that these inscriptions were all made during the same period in the same year (cf. nn. 86 and 109).}\footnote[109]{There may have been good reason for such concern. Following the calendrical reconstructions of 董 (1945), pt. 2, ch. 9, pp. 4a-8a, which would place the 拼编 12-21 divinations in the twelfth month of 吳丁的 twenty-ninth year, with the day hsin-yu being equivalent to 24 November 1311 B.C. (see 董 [1945], pt. 2, \ref{ch:1} (see ch. 1), p. 25b), 长 拼'üan concludes that in the period from the third to the twelfth month of that year seven statelets were at war with the 商 (拼编 12, Tse-tsung [1968], pp. 174-175). But given the k'ao-shih, p. 35; see too, Chow lack of numbered months in all but two of the inscriptions he cites, and the lack of year numbers in all, it is uncertain that these events fell within the same year and in the months he specifies. And it is far less certain that 董's absolute dates are reliable; I myself would, at the moment, prefer dates of about 1198 to 1180 B.C. (see appendix 4, sec. 6). The questionable validity of 董's chronological reconstructions will be discussed in Studies.} 105
the full significance of charges 3 to 6 has to be obtained not just from the inscriptions in
the set but also from those on related scapulas and plastrons; other sets should presumably
be read with the same possibility in mind.) The diviner presented the topics to the spirits
as pairs of positive-negative, complementary charges. In effect the charges offered two
alternatives; if the king followed a certain general to attack a certain tribe he either would
or would not receive assistance. The assistance was presumably to come from Ti, the high
god; that Ti was not specifically mentioned is another instance of routine abbreviation.106
For each of these first six charges, three positive and three negative, one crack was
made on each shell in the sequence indicated by the crack numbers. Thus, the topic of
Wang 成 and the attack on the 夏 韦, proposed in the form of complementary
charges (sentences 1 and 2), was cracked twice on 拼编 12, twice on 拼编 14, and
so on for a total of ten crackings in all. 107 The topic of 知 郭 and the attack on the
八 was also cracked ten times on the five shells (charges 3 and 4); and the divination
was then repeated ten more times-charges 5 and 6 virtually duplicating, though with a
different emphasis, charges 3 and 4-for a total of twenty cracks in all. In addition to the
thirty cracks made for these first six charges, we can discover at least 130 more cracks made
on other bones and shells in a two-day period, probably about this same question of
strategy. 108 Such pyromantic attention suggests that Wu Ting was deeply concerned by
the topic (and also, perhaps, that a satisfactory response was not easy to obtain). 109
% source: scan 102, printed 83
CCCC D

We are uncertain about the outcome. Prognostications and verifications were never
recorded for divinations about following particular officers to attack particular enemies.
As for crack notations, none is recorded on our set for these six charges, though it is possible
that notations might have been carved on the missing fragments.\footnote[110]{Crack notations might have been recorded on missing fragments for 拼编 12.2; 12.5; 12.6; 14.3; 16.1; 18.2; 18.4; 18.6; 20.2; 20.5. III. 拼编 24 + 京 1266 (拼编 24, k'ao-shih, pp. 48-50), probably divined on the same day; the crack notation shang-chi appears by charges 7 and 10.} 110 The fact that two
``highly auspicious'' crack notations appear by the ``follow 知 郭 to attack the 八''
charges on a related plastron, however, suggests that the king may have decided to join
that campaign rather than follow Wang 成 against the 夏 韦.111
The remaining four charges on the front of the shell (inscriptions 7 to 10), written in
smaller calligraphy, concern a toothache, presumably the king's.\footnote[112]{Only Wu Ting (钱-编 4.4.2; 6.33.1) and his consort Fu Hao (易编 3164) are recorded as suffering from toothache (S301.3-4); in all other cases, the subject-presumably the kingis omitted.} 112 None of these inscriptions record a date or diviner's name, and no dated inscriptions on other bones or shells
may be related conclusively to these. On the basis of the dates recorded in the prognostications (see table 12), these charges were probably divined some ten days after the campaign
divinations.\footnote[113]{It is reasonable to suppose that when calligraphy of different sizes appears on the same plastron, the inscriptions written in larger calligraphy record divinations that were made first. On 拼编 1, for example, the small and medium inscriptions appear to have been carved after the large inscriptions and squeezed in among them (on the meaning of these descriptive terms, see sec. 4.3.1.3). The cyclical dates support such a hypothesis; divinations 1 to 4, written in large, bold calligraphy (* tieng height, 9.5-15 mm.), were divined on jen-tzu (day 49) and kuei-ch'ou (day 50). All the other divinations (*tieng height, 5.5-7 mm.) were made later: on keng-shen (day 57), hsin-yu (day 58), kuei-hai (day 60), yi-ch'ou (day 2), and ping-yin (day 3). Further, the subject matter of the small inscriptions is supplementary to that of the large. When the subject matters are unrelated, however, the repetition of the kan-chih cycle may make firm conclusions impossible. E.g., was 拼编 207.3 (large calligraphy; divined on day 33) inscribed before 拼编 207.1 (small calligraphy; divined on day 24)? If so, the front of the plastron would have been used for divination for a period of fifty-one days. If 207.1 was divined before 207.3, however, an interval of only nineteen days would have been involved (the period from day 24 to the day 43 of 207.4).} 113
Charge 7 is not paired, at least not obviously. It states the intention of yu sacrificing a
dog to Fu Keng and mao sacrificing a sheep. Charge 8, whose precise meaning is uncertain,
appears to pray that a sick tooth be removed and that the ting sacrifice be favorable. The
placement of these two inscriptions, balanced on either side of the central axis, suggests
that, though not complementary in content, they formed a pair of linked divinations, the
yu and mao sacrifices being combined with the prayer to lead away the sick tooth.\footnote[114]{This conclusion is strengthened by the link between the yu, mao, and ting sacrifices recorded at 拼编 340.2 and 340.5. 85 U} 114 No
crack notations appear for any of the five cracks associated with charge 7 on the five
plastrons. 小迟 appears by the number 2 crack of charge 8 on 拼编 14, indicating
that the offering of prayer and ting sacrifice and-if one links charges 7 and 8 togetherthe sacrifices to Fu Keng would be ``slightly auspicious.''
Charges 9 and 10 are a complementary, positive-negative pair, divining whether or
% source: scan 103, printed 83
not the sick tooth would prove favorable. No crack notations appear beside the ten cracks
made for this pair of charges (though it is possible that a crack notation might have been
recorded on the missing fragments of 拼编 12).

\subsubsection[The Back of the Plastron]{\origsecnum{3.7.4.2}The Back of the Plastron}
\label{sec:chapters-ch03:3.7.4.2}

The naming of Fu Keng in charge 7 on the front suggests that it and the related
toothache inscriptions charges 8 to 10 on the front) may be linked to the charges on the
back of the plastrons, inscriptions 1 to 8 on 拼编 13, 15, 17, 19, and 21. These record
four complementary charge pairs of the form ``It is because of Fu 甲/ It is not because of
Fu 甲 and the like. Since we know from other inscriptions that toothaches were caused
by ancestral curses,\footnote[115]{E.g., 易编 4511; 4600; 4626, etc. (S301.4). Cf. 长 宗 (1970), pp. 34-45. crack notations, are missing: for Fu 甲, H8 on 拼编 13, H3 on 拼编 17, H4 on 拼编 19; for Fu Keng, HII on 拼编 13, H12 on} ,115 it is reasonable to suppose that these charges on the back were
attempts to determine which ancestor was causing the toothache recorded on the front.
I refer to such charges, which are part of a larger enquiry and which often, as in this case,
lack any preface, as ``subcharges.''
The ancestors to be identified in this instance were all 吳丁的 classificatory fathers
--known to Ssu-ma 钱 as Yang 甲 (K17), 盘 Keng (K18), Hsiao 新 (K19),
and Hsiao 乙 (K20)—and divined here in the order in which they ascended the throne.\footnote[116]{See table 1. On the identification of these 拼编 19; for Fu 新, H13 and H15 on ancestors, \ref{sec:chapters-ch04:4.3.1} (see see sec. 4.3.1).I.}116
As we have speculated\ref{sec:chapters-ch03:3.7.2} (see  sec. 3.7.2) and table 11), cracks about Fu 甲, Fu 新, and Fu
乙 were made twice for the positive and twice for the negative subcharge on each of the
five plastrons; cracks about Fu Keng were made once for the positive and once for the
negative subcharge on each of the five plastrons. The fact that Fu Keng was chosen to be
the recipient of the sacrifice recorded in charge 7 on the front of the shells suggests that
the diviners had identified him as the ancestor responsible for the toothache; if so, the
subcharges recorded on the back of the shell must have been cracked before charge 7 on
the front. The crack notations, ``highly auspicious'' and ``slightly auspicious'' associated
with the Fu 新 subcharges (sec. 3.7.3) suggest that Fu 新 may have been absolved of
responsibility for the toothache. Fu 甲 and Fu 乙 were presumably absolved in the
same way, perhaps by auspicious crack notations on fragments now lost, or by inconclusive
divinations.\footnote[117]{Using the hollow numbers given in fig. 23 and table 12 the following hollows and their fragments, which might have borne the relevant 拼编 13, H14 on 拼编 19; for Fu 乙, H19, H20, H23, and H24 on 拼编 13 and 19, H20, H23, and H24 on 拼编 21.} 117 There are no crack notations associated with subcharges 3 and 4 on the
back of the shell (see table 11) to tell us why Fu Keng was the ancestor identified, but
that to be expected since inauspicious crack notations were never recorded (sec. 2.6).
The subcharges for Fu Keng may have been cracked fewer times than those for his brothers
% source: scan 104, printed 84
because the diviners had identified his inauspicious influence after just two crackings,
making further divination in his case unnecessary.
If this reconstructed scenario is correct, subcharges 1 to 8 carved on the back of the
plastrons record 吳丁的 attempt to discover which classificatory father was causing
his toothache. Charges 7 to 10 on the front indicate that 盘 Keng had been identified
as the uncle responsible and that divinations were then made to determine if offering a
dog and a sheep, while praying to remove the sick tooth, would have a favorable outcome.
The response appears to have been ``slightly auspicious,'' appropriately cautious, perhaps,
in the case of toothache, where no miracle cure could be expected.
Finally, we must attempt to make sense of the two prognostications, 拼编 17.9
and 19.9 (my numbering). The thickness of the carved strokes and the large size of the
graphs may be explained partly by the roughness of the plastron back and partly, perhaps,
by the desire to ``display'' the king's prognostications,\footnote[118]{For display inscriptions, see \ref{ch:2} (see ch. 2), n. 90. The inscriptions in our set are not true display inscriptions, by my definition, for the prognostication is not linked to a charge or verification.} 118 which, in this case, forecast
whether the natural phenomenon would be auspicious or inauspicious according to the
day on which it occurred. The pattern, ``If it be a wu day when there is a natural
phenomenon, etc.,'' indicates that the prognostications were a response to a series of
subcharges which, though unrecorded on our set, must have been of the form “A
natural phenomenon will occur on such-and-such a day,'' each crack indicating whether
the phenomenon was to be considered auspicious or not for the particular day associated
with that crack. These unrecorded subcharges may tentatively be assigned to the
inscriptionless cracks identified in sec.\footnote[119]{If the prognostications were associated with the eight inscriptionless cracks in H25 to H32 (sec. 3.7.2), it is possible to imagine, as in table 12, a series of charges about the natural phenomenon occurring on eight days in sequence. Such an interpretation would require that an auspicious crack notation be found by one of the eight hollows, preferably H29. No crack notations appear, in fact, by any of the hollows, but that may be because four of the five plastrons (all except 拼编 16) lack the bottom halves (or more) of the hypoplastrons and all of the xiphiplastrons, which would have contained all but one (H25 is preserved in 拼编 19) of the missing hollows, H25 to H32. That is to say, of the forty relevant hollows, we now possess only nine; of the five H29s, we possess only one. The hypothesis is well supported by the fact that the sequence of days given in table 12 accords with the sequence of days recorded in the prognostications. We cannot be certain that the sequence is correct since three of the four dates recorded in the prognostications give only the kan-stem and lack the chihbranch. But if this sequence is correct, table 12 reveals that the king, in divining about the sacrifice days when the phenomenon would occur, started by proposing more distant days (when, perhaps, the toothache was more likely to have taken its course and departed) and then proposed days progressively closer, the last day, hsin-wei (?) (day 8), being cracked in the last hollow H32) on the plastron. If this was the day of divination, the plastron set was used for the eleven-day period from day 58 to day 8. Such an interpretation is supported by the fact that another plastron set spanned a similar eleven-day period, from day 42 to day 52 (拼编 35.1; 35.5). A study of 拼编 235.1-2 (fig. 11) and of 拼编 247.1-2 (fig. 12) plus 248.7 (fig. 13) will support the view that individual cracks were associated with individual days in this way. Further evidence for the existence of the putative subcharges will be presented in Studies. 87 U U ]} 3.7.2.119
It is puzzling that the prognostication topic, the natural phenomenon, is mentioned
% source: scan 105, printed 85
not
NKKK

乙六六奋发\textasciitilde{}早人出吸于进Knot (1\textasciitilde{}HED DOPE \listitem{31}
in none of the charges on the shell. Given the absence of inscriptions which would provide
firm corroboration, the interpretation which follows is tentative. A study of another
inscription in which the natural phenomenon appears suggests that it was associated
in the 商 diviner's mind with ancestral sacrifices:
\inscriptionsection{INSCRIPTION}

\begin{inscriptionblock}
\inscriptionref{拼编 207.3 + 208.2}

(PREFACE:) 裂 on ping-shen (day 33), Ch'üeh divined:
(CHARGE:) ``On the following yi-ssu (day 42), offer wine libation to 夏 乙.''
(PROGNOSTICATION:) The king, reading the cracks, said: ``When we offer wine libation
there will be harm. There will perhaps be a natural phenomenon.''
(VERIFICATION:) On yi-ssu (day 42), we offered wine libation. In the early morning
it rained; at the fa sacrifice it stopped raining; at the fa sacrifice
to 仙 (=成 [?], K1), it also rained. We offered shih and
mao sacrifice to the Bird Star(s). (Plastron back:) On the evening
of yi-ssu (day 42) there was a natural phenomenon in the
west.\footnote[120]{拼编 207.3 + 208.2 (part 2); originally published as 易编 6664 (S94.2).}
\end{inscriptionblock}

An analogous situation may be proposed for our plastron set. The response to the charge,
``Yu sacrifice a dog to Fu Keng (and) mao sacrifice a sheep'' (charge 7 on the front) would
have been a prognostication starting, ``When we offer yu and mao sacrifice there will be
harm....'' This was unrecorded, and only the second, more specific part of the prognostication appears on our plastron set: ``There will perhaps be a natural phenomenon.
If it be a wu day when there is a natural phenomenon, it will be inauspicious, etc.''
The fact that no divination was made to determine the day of the sacrifice suggests that
it was to be offered on a daily basis-likely enough in case of toothache-and that the
efficacy of the rituals was to be judged in part by the conjunction, sometimes auspicious,
sometimes not, of the sacrifice day with the natural phenomenon. One may speculate that
the sacrifice charge, together with the prognostications (拼编 17.9; 19.9), demonstrated
that the king, whose campaign plans may have been delayed by his toothache, had not
lost control of the situation. The inscriptions recorded that he was sacrificing, that the
spirits were expected to respond (with thunder or lightning, perhaps), and that he had
already divined the meaning of their response.

\subsection[Making the Record]{\origsecnum{3.7.5}Making the Record}
\label{sec:chapters-ch03:3.7.5}

At the end, a record was made; the engraver carved the crack numbers, crack
notations, charges, and prognostications into the front and back of the plastrons. Whether
this was done at one time or whether, for example, crack numbers, crack notations, charges,
and subcharges were carved first and the prognostications only added later we cannot
% source: scan 106, printed 86
tell; nor can we tell whether the campaign inscriptions were carved before the toothache
divinations were even cracked.\footnote[121]{See \ref{ch:2} (see ch. 2), n. 90.}121
Finally, colored pigment (sec. 2.12) was rubbed into the graphs carved on the front
and back of the shells, as shown in table 8. In most cases, the color has faded badly, and
we cannot be sure if all the graphs were originally colored; nor, pending chemical analysis,
can we be sure if the brown-black distinction reflects a difference in the original colors
or, as seems more likely, a different rate of fading, due perhaps to the orientation of the
plastrons in the ground.\footnote[122]{Cf. sec. 2.12.1. According to the k'ao-shih, p. 40, the graphs of 拼编 17.1-2 seem not to have been colored. The comment on p. 36 of the k'ao-shih that 拼编 13.7 and 13.8 on the bridge were not colored is a typographical error; it was the marginal notation, 拼编 13.9, that was not colored (刘 Yüan-lin, letter of 23 September 1975, citing 长 拼'üan). Further analysis concerning the fate of this set after it was placed in the ground will be presented in Studies. 89 U} 122

\subsection[Conclusions]{\origsecnum{3.7.6}Conclusions}
\label{sec:chapters-ch03:3.7.6}

In sum, this set of five plastrons was used first to divine 吳丁的 campaign
strategy for the season in question (thirty crackings), and it was then reused, about ten
days later, to determine which ancestor was causing his toothache (seventy crackings),
whether or not sacrifice to Fu Keng would be efficacious (twenty crackings), and which
sacrifice days would and which would not be auspicious (forty crackings).
It is worth stressing the time that must have been involved. If the diviners made one
crack a minute (and that seems an overly rapid estimate), it would have taken over two
hours to perform the toothache and sacrifice-day divinations in this set; if they made one
crack every two minutes, it would have taken over four hours, and so on. It is possible, of
course, that cracks were made simultaneously, all the number I cracks on a plastron (or
even on all five plastrons) being burned at the same time for the same topic. Such a procedure would have reduced the time involved, but not the amount of labor. If we include
the estimated ten man-hours required to clean, saw, and polish the five plastrons, the
eighty man-hours required to prepare the hollows (sec. 3.7.4), the additional thirty minutes
(at a minimum) required to crack the campaign divinations, and the five hours (at least)
needed to carve and color the inscriptions, we can see that divining this particular plastron
set would have required some one hundred man-hours of specialized labor, or twenty
hours per plastron. If we assume that a scapula would have involved similar effort and
apply that twenty-hour figure to the total estimated corpus of 69,000 scapulas and 69,000
plastrons (appendix 3, sec. 4), we may estimate that some three million man-hours were
devoted to pyromancy during the historical period; or, put another way, to divine 1.3
scapulas and 1.3 plastrons every day (appendix 3, sec. 4) would have required about fifty
man-hours a day. Crude though these estimates are--and they are probably conservative―
they indicate, once again, the importance the 商 attached to the divination ritual.
% source: scan 107, printed 87
I do not insist upon the accuracy of the above interpretation in every detail. The
meaning of some of the charges and terms used is still obscure; and the relation between
the charges and the cracks remains hypothetical. But this kind of attempt to read an
inscription in full context and to discover its rationale must be undertaken, whenever
possible, before an inscription is used as a historical source. Such close analysis, with one's
nose to the bone as it were, may appear tedious at times-it is undoubtedly tedious to
read about it in this way—but it is the kind of analysis that will, on occasion, provide
insights that advance our knowledge of 商 divination and culture in fresh and
unexpected ways. Scholars must understand the other inscriptions on the same bone
or shell, both front and back, and the possible relations between them; they must look
for crack notations; they must, using crack numbers, look for other inscriptions from the
same set; and they must, using 岛 Sōrui, look for related inscriptions. They must try
to discover what was going on-on the oracle bone, in the 商 court, in the diviner's
mind.
